\documentclass[11pt,a4paper]{article}

\usepackage[utf8]{inputenc}
\usepackage[T1]{fontenc}
\usepackage{lmodern}
\usepackage{textcomp}
\usepackage[margin=2.2cm]{geometry}
\usepackage{amsmath}
\usepackage{amssymb}
\usepackage{booktabs}
\usepackage{longtable}
\usepackage{array}
\usepackage{xcolor}
\usepackage{fancyvrb}
\usepackage{framed}
\usepackage{hyperref}

\definecolor{shadecolor}{RGB}{245,246,248}
\definecolor{num}{RGB}{30,90,160}
\definecolor{alerte}{RGB}{170,40,30}
\hypersetup{colorlinks=true,linkcolor=num,urlcolor=num}

\setlength{\parindent}{0pt}
\setlength{\parskip}{5pt}
\sloppy

% Libelle "L. 115" devant chaque explication de ligne
\newcommand{\ligne}[1]{\item[\textcolor{num}{L.\,#1}]}
\newenvironment{explication}
  {\begin{description}\setlength{\itemsep}{3pt}\setlength{\parskip}{2pt}}
  {\end{description}}
\newcommand{\attention}{\par\textcolor{alerte}{\textbf{Attention :}}\ }
\newcommand{\fichier}[1]{\texttt{#1}}

\title{Le RAG de Ghabetna\\[4pt]
\large Enchaînement des processus et explication ligne par ligne de \texttt{rag/chunker.py}}
\author{\texttt{backend/chatbot\_ms}}
\date{}

\begin{document}
\maketitle
\tableofcontents
\newpage


% ============================================================================
\section{Vue d'ensemble : les deux processus du RAG}
% ============================================================================

Le RAG de Ghabetna n'est pas une seule chaîne de fichiers exécutés l'un après l'autre. Il est
composé de \textbf{deux processus distincts} qui ne s'exécutent pas au même moment et ne
passent pas par les mêmes fichiers :

\begin{enumerate}
  \item \textbf{L'indexation} (hors ligne) : exécutée une seule fois, puis à chaque
        modification du corpus. Elle transforme les fichiers \fichier{.md} en index
        consultables.
  \item \textbf{La réponse à une question} (en ligne) : exécutée à chaque requête d'un
        utilisateur. Elle lit les index déjà construits.
\end{enumerate}

Trois précisions importantes sur l'ordre des fichiers :

\begin{itemize}
  \item \fichier{store.py} n'intervient \textbf{pas} pendant l'indexation : il ne construit
        rien, il \emph{lit} les fichiers produits par l'indexation, au moment de répondre.
  \item \fichier{synonymes.py} et \fichier{calendrier.py} n'interviennent \textbf{pas}
        pendant l'indexation : ils travaillent uniquement sur la question de l'utilisateur.
  \item Le cache est \textbf{consulté en premier}, avant toute recherche, et
        \textbf{alimenté en dernier}, après la génération.
\end{itemize}

\fichier{config.py} est utilisé par tous les fichiers : il regroupe tous les réglages
(chemins, modèle, tailles, clés).

% ----------------------------------------------------------------------------
\subsection{Processus A : indexation des données}
% ----------------------------------------------------------------------------

Commande : \texttt{python -m scripts.build\_index}

\begin{shaded}
\begin{Verbatim}[fontsize=\footnotesize]
scripts/build_index.py        (chef d'orchestre : lance 2 processus Python separes)
|
+-- PROCESSUS 1 : scripts/build_embeddings.py           (lourd : charge torch)
|     |
|     +-- ETAPE 1 : rag/chunker.py  -> chunker_corpus()
|     |       lit   data/corpus/<domaine>/*.md  +  _manifest.json
|     |       ecrit data/index/chunks.json                (~320 chunks)
|     |
|     +-- ETAPE 2 : rag/indexer.py  -> embedder(chunks)
|             charge le modele bge-m3, transforme chaque chunk en vecteur
|             ecrit data/index/vecteurs.npy               (matrice 320 x 1024)
|                   data/index/vecteurs.meta.json         (empreinte du cache)
|
+-- PROCESSUS 2 : scripts/build_stores.py               (leger : sans torch)
      |
      +-- ETAPE 3 : rag/indexer.py  -> indexer_bm25(chunks)
              ecrit data/index/bm25.pkl                   (index lexical)
\end{Verbatim}
\end{shaded}

\textbf{Pourquoi deux processus ?} L'embedding est l'étape lente (environ 11 minutes sur
CPU). Comme son résultat, \fichier{vecteurs.npy}, est enregistré sur disque, l'étape BM25
peut être relancée en quelques secondes sans recharger le modèle. Séparer les processus
évite aussi les conflits de bibliothèques natives (DLL) entre \texttt{torch} et les autres
bibliothèques.

\textbf{Résultat :} trois fichiers dans \fichier{data/index/}.

\begin{longtable}{>{\ttfamily}l p{9.5cm}}
\toprule
\normalfont\textbf{Fichier} & \textbf{Contenu} \\
\midrule
chunks.json  & la liste des chunks : textes et métadonnées \\
vecteurs.npy & un vecteur de 1024 nombres par chunk, dans le même ordre ; c'est le
               « magasin vectoriel » lu par \fichier{store.py} \\
bm25.pkl     & l'index lexical BM25 et la liste des identifiants dans le même ordre \\
\bottomrule
\end{longtable}

% ----------------------------------------------------------------------------
\subsection{Processus B : traitement d'une question utilisateur}
% ----------------------------------------------------------------------------

\textbf{Au démarrage du service} (fonction \texttt{lifespan} de \fichier{app/main.py}),
une seule fois :

\begin{shaded}
\begin{Verbatim}[fontsize=\footnotesize]
Retriever()  -> charge le modele bge-m3 (10 a 20 s)
             -> NumpyStore() lit vecteurs.npy + chunks.json      (store.py)
             -> lit bm25.pkl
\end{Verbatim}
\end{shaded}

\textbf{Puis à chaque appel de} \texttt{POST /api/chat/} :

\begin{shaded}
\begin{Verbatim}[fontsize=\footnotesize]
 Flutter envoie { question, historique, specialite }
      |
 (1) config.domaines_pour("chasseur")          -> ["chasse", "app"]
      |
 (2) calendrier.detecter_date(question)        -> la question contient-elle une date ?
      |
 (3) cache.lire()     (seulement si : pas d'historique ET pas de date)
      |-- trouve  -> reponse immediate, FIN
      |
 (4) retriever.rechercher(question, domaines, question_precedente)
      |    +-- _dense() : bge-m3 encode la question -> store.query()     (le SENS)
      |    +-- _bm25()  : synonymes.enrichir(question) -> tokeniser -> BM25   (les MOTS)
      |    +-- _rrf()   : fusion des deux classements -> top 8 chunks
      |
 (5) calendrier.analyser_question_struct(question)
      |    espece + date ?  -> verdict calcule en Python (autorise / refuse / passee)
      |
 (6) generator.generer(question, chunks, historique, specialite, verdict)
      |    prompt systeme + [VERDICT] + [EXTRAITS] + QUESTION  -> Groq (LLM)
      |
 (7) cache.ecrire()   (memes conditions qu'en (3))
      |
 (8) reponse JSON : { reponse, verdict, sources, domaines, depuis_cache, duree_ms }
\end{Verbatim}
\end{shaded}

\textbf{Où interviennent synonymes et calendrier ?}
\begin{itemize}
  \item \fichier{synonymes.py} est appelé \emph{à l'intérieur} du retriever, et
        \textbf{uniquement} du côté BM25 : on ajoute à la question les mots du vocabulaire
        officiel (« prix » $\rightarrow$ « redevance, montant, dinars »). On ne le fait
        pas côté dense, car ajouter des mots déplacerait le vecteur de la question vers une
        moyenne floue.
  \item \fichier{calendrier.py} est appelé \emph{à côté} du retriever, pas dedans. Il sert
        deux fois : à l'étape (2) pour décider si on utilise le cache, et à l'étape (5)
        pour calculer le verdict.
\end{itemize}

\textbf{Pourquoi le cache est-il désactivé avec un historique ou une date ?} Avec un
historique, la même question peut appeler une réponse différente selon ce qui précède
(« et pour le lièvre ? »). Avec une date relative, « demain » ne désigne pas le même jour
d'un jour à l'autre. Dans ces deux cas, une réponse en cache serait fausse.

\newpage
% ============================================================================
\section{Explication ligne par ligne de \texorpdfstring{\texttt{rag/chunker.py}}{rag/chunker.py}}
% ============================================================================

\subsection*{Rôle du fichier}

\fichier{chunker.py} est la \textbf{première étape} de l'indexation. Il lit les fichiers
Markdown de \fichier{data/corpus/<domaine>/} et les découpe en petits morceaux autonomes
appelés \emph{chunks}, enregistrés dans \fichier{data/index/chunks.json}.

Sa règle fondamentale est : \textbf{un chunk = un article de loi}. Le découpage est
\emph{sémantique} : on cherche le motif « Article N » dans les titres, au lieu de se fier au
niveau \texttt{\#\#} ou \texttt{\#\#\#}, qui varie d'un fichier à l'autre.

\subsection*{Exemple de chunk réellement produit}

Extrait de \fichier{chunks.json} (première ligne du tableau de l'article premier de l'arrêté) :

\begin{shaded}
\begin{Verbatim}[fontsize=\footnotesize]
{
 "id": "arrete_2025_art1_lievre_perdrix_caille_sedentai",
 "texte_embedding": "Arrêté chasse 2025/2026, Article premier. Lièvre, perdrix, ...
                     Lièvre, perdrix, ... Date d'ouverture : 5 octobre 2025, ...",
 "texte_indexe":    "Arrêté du 27 août 2025 (JORT n° 107 du 29 août 2025) — Titre
                     premier : Réglementation générale — Article premier — Lièvre, ...",
 "texte_affiche":   "Lièvre, perdrix, ... Date d'ouverture : 5 octobre 2025,
                     Date de fermeture : 7 décembre 2025.",
 "metadata": { "domaine": "chasse", "source_id": "arrete_2025", "article": "1",
               "type": "tableau", "validite": "saison_2025_2026", ... }
}
\end{Verbatim}
\end{shaded}

\subsection*{Comment lire la suite}

Le code est reproduit \textbf{intégralement}, par blocs, avec les numéros de ligne réels
du fichier. Sous chaque bloc, chaque ligne est expliquée (\textcolor{num}{L.\,n}). Les lignes
vides et les commentaires sont aussi signalés, pour qu'aucune ligne ne soit sautée.

\subsection{En-tête, docstring et imports (lignes 1 à 30)}
\begin{shaded}
\begin{Verbatim}[numbers=left,firstnumber=1,numbersep=6pt,fontsize=\footnotesize,xleftmargin=1.2em]
#!/usr/bin/env python3
# -*- coding: utf-8 -*-
\end{Verbatim}
\end{shaded}

\begin{explication}
\ligne{1} \emph{Shebang}. Sous Linux ou macOS, si on exécute le fichier directement
(\Verb@./chunker.py@), le système lit cette ligne pour savoir quel interpréteur utiliser.
\Verb@/usr/bin/env python3@ cherche \Verb@python3@ dans le \texttt{PATH}. Cette ligne est
sans effet sous Windows et quand on lance le module avec \Verb@python -m@.
\ligne{2} Déclaration d'encodage (norme PEP~263). En Python~3, l'UTF-8 est déjà l'encodage par
défaut : la ligne est gardée par convention. Elle rappelle que le fichier contient des
caractères accentués (é, «, »).
\end{explication}

\begin{shaded}
\begin{Verbatim}[numbers=left,firstnumber=3,numbersep=6pt,fontsize=\footnotesize,xleftmargin=1.2em]
"""
rag/chunker.py — ETAPE 1 : decoupage des .md en chunks
=======================================================

REGLE FONDAMENTALE :
    Un chunk = un ARTICLE de loi.
    Le decoupage est SEMANTIQUE (on cherche le motif "Article N" dans le
    titre) et non SYNTAXIQUE : on ne peut pas se baser sur le niveau ## ou
    ###, car il varie d'un fichier a l'autre. Dans l'arrete les articles
    sont en H2 ; dans le code forestier ils sont en H2 pour le Titre I et
    en H3 pour le Titre II.

MULTI-DOMAINES :
    Le nom du sous-dossier de data/corpus/ DEVIENT la metadonnee "domaine".
    Ajouter l'apiculture = creer data/corpus/apiculture/ et y deposer des .md.
    Aucune ligne de code a modifier.

Ce module ne depend d'aucune librairie externe : il se teste seul.
"""
\end{Verbatim}
\end{shaded}

\begin{explication}
\ligne{3} Ouverture d'une chaîne entre triples guillemets. Placée en tout début de module,
elle devient la \emph{docstring} du module : ce n'est pas du code exécuté mais une
documentation, accessible par \Verb@chunker.__doc__@ ou \Verb@help(chunker)@.
\ligne{4--5} Titre du module et soulignement décoratif.
\ligne{6} Ligne vide de mise en forme.
\ligne{7--13} La règle fondamentale : un chunk = un article. On repère les articles par le motif
« Article N » et non par le niveau de titre, car dans l'arrêté les articles sont en
\texttt{\#\#} (H2), alors que dans le code forestier ils sont en H2 pour le Titre~I et en H3
pour le Titre~II.
\ligne{14} Ligne vide.
\ligne{15--18} Principe multi-domaines : le nom du sous-dossier de \fichier{data/corpus/}
devient la métadonnée \texttt{domaine}. Ajouter un domaine consiste à créer un dossier et à y
déposer des \fichier{.md}, sans modifier le code.
\ligne{19} Ligne vide.
\ligne{20} Le module n'utilise que la bibliothèque standard : on peut le tester sans
installer \texttt{torch} ni \texttt{sentence-transformers}.
\ligne{21} Fermeture de la docstring.
\end{explication}

\begin{shaded}
\begin{Verbatim}[numbers=left,firstnumber=22,numbersep=6pt,fontsize=\footnotesize,xleftmargin=1.2em]

from __future__ import annotations

import json
import re
import unicodedata
from pathlib import Path

from . import config
\end{Verbatim}
\end{shaded}

\begin{explication}
\ligne{22} Ligne vide.
\ligne{23} \Verb@from __future__ import annotations@ (PEP~563) : les annotations de type ne sont
plus évaluées à l'exécution mais conservées sous forme de texte. Cela permet d'écrire
\Verb@str | None@ ou \Verb@dict[int, str]@ même sur une version de Python antérieure à 3.10.
Cet import doit être la première instruction du fichier, juste après la docstring.
\ligne{24} Ligne vide.
\ligne{25} \Verb@json@ : lire les \fichier{\_manifest.json} et écrire \fichier{chunks.json}.
\ligne{26} \Verb@re@ : expressions régulières, pour reconnaître titres, articles et tableaux.
\ligne{27} \Verb@unicodedata@ : base de données Unicode, utilisée pour retirer les accents.
\ligne{28} \Verb@Path@ : manipulation de chemins de fichiers orientée objet (module
\Verb@pathlib@).
\ligne{29} Ligne vide.
\ligne{30} \Verb@from . import config@ : \emph{import relatif}. Le point désigne le paquet
courant, \Verb@rag@. On récupère ainsi \Verb@config.TAILLE_MAX_CHUNK@,
\Verb@config.CORPUS_DIR@, etc. Conséquence : le fichier doit être lancé comme module
(\Verb@python -m rag.chunker@) ; un lancement direct (\Verb@python rag/chunker.py@) lève
\emph{attempted relative import with no known parent package}.
\end{explication}

\subsection{Les expressions régulières (lignes 31 à 54)}
\begin{shaded}
\begin{Verbatim}[numbers=left,firstnumber=31,numbersep=6pt,fontsize=\footnotesize,xleftmargin=1.2em]

# ============================================================================
# EXPRESSIONS REGULIERES
# ============================================================================

RE_TITRE = re.compile(r"^(#{1,6})\s+(.+?)\s*#*$")
\end{Verbatim}
\end{shaded}

Toutes les expressions régulières sont \emph{compilées} une seule fois, au chargement du
module, avec \Verb@re.compile@. Elles sont ensuite réutilisées sans être réanalysées. Le
préfixe \Verb@r"..."@ (\emph{raw string}) empêche Python d'interpréter les \Verb@\@ : ainsi
\Verb@\s@ arrive tel quel au moteur d'expressions régulières.

\begin{explication}
\ligne{31} Ligne vide.
\ligne{32--34} Bannière de commentaire qui délimite la section.
\ligne{35} Ligne vide.
\ligne{36} \Verb@RE_TITRE@ reconnaît un titre Markdown. Décomposition :
  \begin{itemize}
    \item \Verb@^@ : début de ligne ;
    \item \Verb@(#{1,6})@ : \textbf{groupe 1}, de 1 à 6 dièses ; son nombre de caractères
          donnera le niveau du titre ;
    \item \Verb@\s+@ : au moins un espace, donc \Verb@#hashtag@ n'est pas un titre ;
    \item \Verb@(.+?)@ : \textbf{groupe 2}, le texte du titre ; le \Verb@?@ rend le
          quantificateur \emph{paresseux} (il prend le moins de caractères possible) ;
    \item \Verb@\s*#*@ : espaces et dièses fermants facultatifs (syntaxe
          \Verb@## Titre ##@) ; grâce au quantificateur paresseux, ils ne sont pas avalés
          par le groupe 2 ;
    \item \Verb@$@ : fin de ligne.
  \end{itemize}
  Exemple : \Verb@### Article 165@ donne groupe~1 = \Verb@###@ et groupe~2 =
  \Verb@Article 165@.
\end{explication}

\begin{shaded}
\begin{Verbatim}[numbers=left,firstnumber=37,numbersep=6pt,fontsize=\footnotesize,xleftmargin=1.2em]

# Detecte un ARTICLE quel que soit son niveau de titre.
# Gere "Article premier", "Article 12", "Article 134 bis".
RE_ARTICLE = re.compile(r"^Article\s+(premier|\d+\s*bis|\d+)\b", re.IGNORECASE)
\end{Verbatim}
\end{shaded}

\begin{explication}
\ligne{37} Ligne vide.
\ligne{38--39} Commentaires : la regex détecte un article quel que soit son niveau de titre.
\ligne{40} \Verb@RE_ARTICLE@ reconnaît un titre d'article :
  \begin{itemize}
    \item \Verb@^Article\s+@ : le mot « Article » en début de chaîne, suivi d'espaces ;
    \item \Verb@(premier|\d+\s*bis|\d+)@ : \textbf{groupe 1}, trois alternatives. \textbf{L'ordre
          compte} : le moteur essaie les alternatives de gauche à droite. Si \Verb@\d+@ était
          placé avant \Verb@\d+\s*bis@, « 134 bis » serait capturé comme « 134 ».
          \Verb@\s*@ accepte « 134 bis » comme « 134bis » ;
    \item \Verb@\b@ : \emph{frontière de mot}. Elle empêche une correspondance partielle
          (« premiers ») ;
    \item \Verb@re.IGNORECASE@ : insensible à la casse (« ARTICLE 5 », « article 5 »).
  \end{itemize}
  Un titre comme « Article 75 (Modifié par la loi n° 2005-13) » est reconnu, car seul le
  début est testé.
  \attention « Article 1er » n'est pas reconnu : entre « 1 » et « er » il n'y a pas de
  frontière de mot, donc \Verb@\b@ échoue. Le corpus actuel écrit « Article premier ».\par
\end{explication}

\begin{shaded}
\begin{Verbatim}[numbers=left,firstnumber=41,numbersep=6pt,fontsize=\footnotesize,xleftmargin=1.2em]

# Sous-section geographique de l'article 12 de l'arrete
RE_GOUVERNORAT = re.compile(r"^Gouvernorat\s+(?:de\s+l'|du\s+|de\s+|d')?(.+)$", re.I)
\end{Verbatim}
\end{shaded}

\begin{explication}
\ligne{41} Ligne vide.
\ligne{42} Commentaire : l'article 12 de l'arrêté est divisé en sous-sections par gouvernorat.
\ligne{43} \Verb@RE_GOUVERNORAT@ :
  \begin{itemize}
    \item \Verb@^Gouvernorat\s+@ : le mot « Gouvernorat » suivi d'espaces ;
    \item \Verb@(?:de\s+l'|du\s+|de\s+|d')?@ : groupe \emph{non capturant} (\Verb@(?:...)@),
          facultatif (\Verb@?@), qui absorbe l'article : « de l' », « du », « de », « d' ».
          Ici aussi l'ordre compte : \Verb@de\s+l'@ doit passer avant \Verb@de\s+@, sinon
          « de l'Ariana » donnerait « l'Ariana » ;
    \item \Verb@(.+)$@ : \textbf{groupe 1}, le nom du gouvernorat jusqu'à la fin ;
    \item \Verb@re.I@ : abréviation de \Verb@re.IGNORECASE@.
  \end{itemize}
  Exemples : « Gouvernorat de l'Ariana » $\rightarrow$ « Ariana » ; « Gouvernorat du Kef »
  $\rightarrow$ « Kef » ; « Gouvernorat de Béja » $\rightarrow$ « Béja ».
\end{explication}

\begin{shaded}
\begin{Verbatim}[numbers=left,firstnumber=44,numbersep=6pt,fontsize=\footnotesize,xleftmargin=1.2em]

# Les blocs > ajoutes dans les .md : deux types traites differemment
RE_NOTE_FIDELITE = re.compile(r"\*\*Note de fidélité\.?\*\*|\*\*Périmètre\.?\*\*", re.I)
RE_AVERTISSEMENT = re.compile(r"\*\*Avertissement de concordance\.?\*\*", re.I)
\end{Verbatim}
\end{shaded}

\begin{explication}
\ligne{44} Ligne vide.
\ligne{45} Commentaire : les blocs de citation \Verb@>@ ajoutés dans les \fichier{.md} sont de
deux types, traités différemment.
\ligne{46} \Verb@RE_NOTE_FIDELITE@ : \Verb@\*\*@ correspond à deux astérisques littéraux (le gras
Markdown ; l'astérisque est échappé car c'est un quantificateur en regex). \Verb@\.?@ accepte
un point facultatif ; \Verb@|@ signifie « ou ». La regex reconnaît donc
« **Note de fidélité.** » ou « **Périmètre.** ».
  \attention cette regex n'est \textbf{jamais utilisée} dans le fichier. Tout bloc
  \Verb@>@ qui n'est pas un avertissement est classé « note » (ligne 168), quel que soit
  son contenu.\par
\ligne{47} \Verb@RE_AVERTISSEMENT@ : même construction, pour « **Avertissement de concordance.** ».
Celle-ci est utilisée à la ligne 168.
\end{explication}

\begin{shaded}
\begin{Verbatim}[numbers=left,firstnumber=48,numbersep=6pt,fontsize=\footnotesize,xleftmargin=1.2em]

# Coquille officielle + forme correcte :  *(sic : « chasseurs »)*
RE_SIC = re.compile(r"\*\(sic\s*:?\s*«?\s*([^»)]*?)\s*»?\)\*")
\end{Verbatim}
\end{shaded}

\begin{explication}
\ligne{48} Ligne vide.
\ligne{49} Commentaire : format d'une coquille du texte officiel suivie de sa forme correcte.
\ligne{50} \Verb@RE_SIC@ reconnaît \Verb@*(sic : « chasseurs »)*@ :
  \begin{itemize}
    \item \Verb@\*\(@ : un astérisque puis une parenthèse ouvrante, littéraux ;
    \item \Verb@sic\s*:?\s*@ : le mot « sic », un deux-points facultatif, des espaces ;
    \item \Verb@«?\s*@ : guillemet ouvrant facultatif ;
    \item \Verb@([^»)]*?)@ : \textbf{groupe 1}, la forme correcte : tout caractère sauf
          \Verb@»@ et \Verb@)@ (classe négative \Verb@[^...]@), en mode paresseux ;
    \item \Verb@\s*»?\)\*@ : guillemet fermant facultatif, puis \Verb@)*@ littéraux.
  \end{itemize}
  Exemple réel du code forestier : « qui cour \Verb@*(sic : « qui court »)*@ » donne le
  groupe~1 = « qui court ».
\end{explication}

\begin{shaded}
\begin{Verbatim}[numbers=left,firstnumber=51,numbersep=6pt,fontsize=\footnotesize,xleftmargin=1.2em]

RE_LIGNE_TABLEAU = re.compile(r"^\s*\|")
# Ligne de donnees d'un tableau markdown : | espece | ouverture | fermeture |
RE_LIGNE_SEPARATRICE = re.compile(r"^\s*\|[\s:|-]+\|\s*$")
\end{Verbatim}
\end{shaded}

\begin{explication}
\ligne{51} Ligne vide.
\ligne{52} \Verb@RE_LIGNE_TABLEAU@ = \Verb@^\s*\|@ : une ligne qui commence (après d'éventuels
espaces) par une barre verticale, donc une ligne de tableau Markdown. Le \Verb@|@ est échappé,
car en regex il signifie « ou ».
\ligne{53} Commentaire. \attention il est mal placé : il décrit une « ligne de données »
(ligne 52), alors qu'il précède la regex de la ligne séparatrice.\par
\ligne{54} \Verb@RE_LIGNE_SEPARATRICE@ = \Verb@^\s*\|[\s:|-]+\|\s*$@ : la ligne
\Verb@|---|:---:|@ qui sépare l'entête des données. La classe \Verb@[\s:|-]@ n'accepte que des
espaces, des deux-points, des barres et des tirets (placé en dernier dans la classe, le
\Verb@-@ est littéral). Une ligne contenant une seule lettre ne correspond donc pas.
\end{explication}

\subsection{Utilitaires (lignes 55 à 97)}
\begin{shaded}
\begin{Verbatim}[numbers=left,firstnumber=55,numbersep=6pt,fontsize=\footnotesize,xleftmargin=1.2em]


# ============================================================================
# UTILITAIRES
# ============================================================================

def plier_accents(txt: str) -> str:
    """
    Aïn -> Ain, Béjà -> Beja.
    Sert UNIQUEMENT a fabriquer des identifiants techniques.
    Le texte affiche garde ses accents.
    """
    d = unicodedata.normalize("NFD", txt)
    return "".join(c for c in d if unicodedata.category(c) != "Mn")
\end{Verbatim}
\end{shaded}

\begin{explication}
\ligne{55--56} Deux lignes vides entre sections de module (convention PEP~8).
\ligne{57--59} Bannière « UTILITAIRES ».
\ligne{60} Ligne vide.
\ligne{61} Définition de \Verb@plier_accents@. \Verb@txt: str@ et \Verb@-> str@ sont des
annotations de type : documentation pour le lecteur et l'éditeur, non vérifiées à l'exécution.
\ligne{62--66} Docstring : la fonction transforme « Aïn » en « Ain ». Elle sert
\emph{uniquement} à fabriquer des identifiants ; le texte affiché garde ses accents.
\ligne{67} \Verb@unicodedata.normalize("NFD", txt)@ : forme de normalisation \emph{NFD}
(décomposition canonique). Chaque caractère accentué est séparé en lettre de base et accent
\emph{combinant} : « é » (U+00E9) devient « e » (U+0065) suivi de « \'{} » (U+0301).
\ligne{68} Expression génératrice : on parcourt chaque caractère \Verb@c@ et on ne garde que ceux
dont la catégorie Unicode n'est pas \Verb@"Mn"@ (\emph{Mark, nonspacing}, c'est-à-dire les
accents combinants). \Verb@"".join(...)@ recolle les caractères restants.
Résultat : « Béja » $\rightarrow$ « Beja ». Les ligatures comme « œ » ne sont pas décomposées.
\end{explication}

\begin{shaded}
\begin{Verbatim}[numbers=left,firstnumber=69,numbersep=6pt,fontsize=\footnotesize,xleftmargin=1.2em]


def slug(txt: str, taille: int = 40) -> str:
    """'Sidi Bouzid' -> 'sidi_bouzid'"""
    t = plier_accents(txt.lower())
    return re.sub(r"[^a-z0-9]+", "_", t).strip("_")[:taille]
\end{Verbatim}
\end{shaded}

\begin{explication}
\ligne{69--70} Lignes vides.
\ligne{71} \Verb@slug(txt, taille=40)@ : \Verb@taille@ est un paramètre avec
\emph{valeur par défaut} ; on peut l'omettre à l'appel.
\ligne{72} Docstring sur une ligne : exemple d'entrée et de sortie.
\ligne{73} Passage en minuscules, puis suppression des accents.
\ligne{74} Trois opérations enchaînées :
  \begin{itemize}
    \item \Verb@re.sub(r"[^a-z0-9]+", "_", t)@ remplace toute suite de caractères qui ne sont
          ni lettres ni chiffres (classe négative, \Verb@+@ = une ou plusieurs) par un seul
          \Verb@_@ ;
    \item \Verb@.strip("_")@ retire les \Verb@_@ au début et à la fin ;
    \item \Verb@[:taille]@ : \emph{slicing}, on garde au plus \Verb@taille@ caractères.
  \end{itemize}
  Exemple : « Gouvernorat de l'Ariana » $\rightarrow$ \Verb@gouvernorat_de_l_ariana@.
\end{explication}

\begin{shaded}
\begin{Verbatim}[numbers=left,firstnumber=75,numbersep=6pt,fontsize=\footnotesize,xleftmargin=1.2em]


def numero_article(titre: str) -> str | None:
    """
    'Article premier' -> '1'  |  'Article 134 bis' -> '134bis'
    None si ce n'est pas un article (Chapitre, Visas, Preambule...).
    """
    m = RE_ARTICLE.match(titre)
    if not m:
        return None
    num = m.group(1).lower().replace(" ", "")
    return "1" if num == "premier" else num
\end{Verbatim}
\end{shaded}

\begin{explication}
\ligne{75--76} Lignes vides.
\ligne{77} \Verb@-> str | None@ : la fonction renvoie soit une chaîne, soit \Verb@None@.
\ligne{78--81} Docstring : « Article premier » $\rightarrow$ \Verb@"1"@, « Article 134 bis »
$\rightarrow$ \Verb@"134bis"@, \Verb@None@ pour un titre qui n'est pas un article.
\ligne{82} \Verb@RE_ARTICLE.match(titre)@ : \Verb@match@ ne teste que le \emph{début} de la chaîne
(contrairement à \Verb@search@, qui cherche partout). Renvoie un objet \Verb@Match@ ou \Verb@None@.
\ligne{83} \Verb@if not m@ : \Verb@None@ est « faux » en Python.
\ligne{84} Pas un article (« Chapitre II », « Visas »…) : on renvoie \Verb@None@.
\ligne{85} \Verb@m.group(1)@ est le texte capturé par le groupe 1. \Verb@.lower()@ normalise la
casse (grâce à \Verb@IGNORECASE@, on a pu capturer « Premier » ou « BIS »).
\Verb@.replace(" ", "")@ transforme « 134 bis » en « 134bis ».
\ligne{86} \emph{Expression conditionnelle} (\Verb@A if condition else B@) : « premier » devient
\Verb@"1"@, sinon on renvoie le numéro tel quel.
\end{explication}

\begin{shaded}
\begin{Verbatim}[numbers=left,firstnumber=87,numbersep=6pt,fontsize=\footnotesize,xleftmargin=1.2em]


def extraire_alias_sic(texte: str) -> list[str]:
    """
    Recupere les formes correctes derriere les *(sic : « ... »)*.
    L'article 176 ecrit "chausseurs" : coquille du texte OFFICIEL.
    On garde "chausseurs" a l'affichage (fidelite juridique) mais on ajoute
    "chasseurs" comme alias pour que BM25 puisse le retrouver.
    """
    return [m.group(1).strip() for m in RE_SIC.finditer(texte)
            if m.group(1).strip() and len(m.group(1).strip()) < 60]
\end{Verbatim}
\end{shaded}

\begin{explication}
\ligne{87--88} Lignes vides.
\ligne{89} Définition : reçoit un texte, renvoie une liste de chaînes.
\ligne{90--95} Docstring : l'article 176 écrit « chausseurs », une coquille du texte
\emph{officiel}. On la garde à l'affichage pour rester fidèle au texte juridique, mais on
ajoute « chasseurs » comme alias pour que BM25 retrouve l'article.
\ligne{96--97} \emph{Compréhension de liste} :
  \begin{itemize}
    \item \Verb@RE_SIC.finditer(texte)@ parcourt \emph{toutes} les occurrences, une par une
          (contrairement à \Verb@match@, qui ne teste que le début) ;
    \item \Verb@m.group(1).strip()@ : la forme correcte, sans espaces autour ;
    \item filtre : on garde l'alias s'il n'est pas vide et fait moins de 60 caractères, ce
          qui écarte les captures aberrantes trop longues.
  \end{itemize}
  L'expression \Verb@m.group(1).strip()@ est calculée trois fois par occurrence : c'est un
  peu redondant, mais le coût est négligeable.
\end{explication}

\subsection{\texorpdfstring{\texttt{decouper\_sections}}{decouper\_sections} : lire la structure du document (lignes 98 à 134)}
\begin{shaded}
\begin{Verbatim}[numbers=left,firstnumber=98,numbersep=6pt,fontsize=\footnotesize,xleftmargin=1.2em]


# ============================================================================
# DECOUPAGE EN SECTIONS
# ============================================================================

def decouper_sections(texte: str) -> list[dict]:
    """
    Produit une liste de sections, chacune avec LA PILE DE SES PARENTS.

    La pile est essentielle : "### Article 165" ne dit pas a quel Titre ni a
    quel Chapitre il appartient. Ce sont ses parents qui le disent, et c'est
    cette information qui constituera le prefixe contextuel.
    """
    sections, courante = [], None
    parents: dict[int, str] = {}
\end{Verbatim}
\end{shaded}

C'est la fonction clé du fichier. Elle lit le Markdown ligne par ligne et produit une
liste de sections. Chaque section mémorise \textbf{la pile de ses titres parents}.

\begin{explication}
\ligne{98--99} Lignes vides.
\ligne{100--102} Bannière « DECOUPAGE EN SECTIONS ».
\ligne{103} Ligne vide.
\ligne{104} Définition : reçoit le texte complet d'un fichier et renvoie une liste de
dictionnaires.
\ligne{105--111} Docstring : « \#\#\# Article 165 » ne dit pas à quel Titre ni à quel Chapitre
il appartient. Ce sont ses parents qui le disent, et ils serviront au préfixe contextuel.
\ligne{112} Affectation multiple (\emph{déballage de tuple}) : \Verb@sections = []@ contiendra
le résultat ; \Verb@courante = None@ est la section en cours de remplissage (aucune au départ).
\ligne{113} \Verb@parents@ est un dictionnaire \{niveau : titre\}, par exemple
\Verb@{1: "Titre II", 2: "Chapitre I"}@. L'annotation \Verb@dict[int, str]@ documente les types.
\end{explication}

\begin{shaded}
\begin{Verbatim}[numbers=left,firstnumber=114,numbersep=6pt,fontsize=\footnotesize,xleftmargin=1.2em]

    for ligne in texte.splitlines():
        m = RE_TITRE.match(ligne)
        if m:
            niveau, titre = len(m.group(1)), m.group(2).strip()
            if courante:
                sections.append(courante)
            # On oublie les parents de niveau >= au notre : en arrivant sur
            # "## Chapitre II" on oublie les "###" du chapitre precedent,
            # mais on garde le "# Titre II".
            parents = {k: v for k, v in parents.items() if k < niveau}
            courante = {"niveau": niveau, "titre": titre,
                        "parents": [parents[k] for k in sorted(parents)],
                        "lignes": []}
            parents[niveau] = titre
        elif courante is not None:
            courante["lignes"].append(ligne)
\end{Verbatim}
\end{shaded}

\begin{explication}
\ligne{114} Ligne vide.
\ligne{115} \Verb@texte.splitlines()@ découpe le texte en lignes. Contrairement à
\Verb@split("\n")@, elle gère aussi les fins de ligne Windows \Verb@\r\n@.
\ligne{116} On teste si la ligne est un titre Markdown.
\ligne{117} \Verb@if m:@ : un objet \Verb@Match@ est « vrai », \Verb@None@ est « faux ».
\ligne{118} \Verb@len(m.group(1))@ compte les dièses : \Verb@###@ donne le niveau 3.
\Verb@m.group(2).strip()@ donne le texte du titre sans espaces.
\ligne{119--120} Un nouveau titre commence : la section précédente est terminée et on l'ajoute au
résultat. Au tout premier titre, \Verb@courante@ vaut \Verb@None@ et rien n'est ajouté.
\ligne{121--123} Commentaire qui explique la ligne suivante.
\ligne{124} \emph{Compréhension de dictionnaire} : on ne garde que les parents de niveau
\textbf{strictement inférieur} au titre courant. En arrivant sur « \#\# Chapitre II », on
oublie les « \#\#\# » du chapitre précédent, mais on garde « \# Titre II ».
\ligne{125--127} Création de la nouvelle section :
  \begin{itemize}
    \item \Verb@"niveau"@ et \Verb@"titre"@ : ceux du titre lu ;
    \item \Verb@"parents"@ : \Verb@sorted(parents)@ trie les \emph{clés} (les niveaux), puis on
          prend les titres dans cet ordre : du plus général au plus précis ;
    \item \Verb@"lignes"@ : liste vide, remplie par les lignes suivantes.
  \end{itemize}
  La pile est calculée \emph{avant} la ligne 128 : une section n'est donc jamais son propre parent.
\ligne{128} Le titre courant devient parent des sections qui suivront.
\ligne{129} \Verb@elif courante is not None@ : ce n'est pas un titre et une section est ouverte.
\ligne{130} La ligne brute (non nettoyée) est ajoutée au contenu de la section courante.
\end{explication}
\attention les lignes situées \textbf{avant le premier titre} du fichier ne vont nulle part :
\Verb@courante@ vaut encore \Verb@None@, elles sont perdues sans avertissement.\par

\begin{shaded}
\begin{Verbatim}[numbers=left,firstnumber=131,numbersep=6pt,fontsize=\footnotesize,xleftmargin=1.2em]

    if courante:
        sections.append(courante)
    return sections
\end{Verbatim}
\end{shaded}

\begin{explication}
\ligne{131} Ligne vide.
\ligne{132--133} À la fin de la boucle, la dernière section n'a pas été ajoutée, car aucun titre
suivant n'a déclenché la ligne 120. On l'ajoute ici.
\ligne{134} Renvoi de la liste des sections.
\end{explication}

\textbf{Exemple d'évolution de la pile :}

\begin{longtable}{l l l}
\toprule
\textbf{Ligne lue} & \textbf{\texttt{parents} après L.124} & \textbf{Pile enregistrée} \\
\midrule
\texttt{\# Titre II}          & \{\}                     & [ ] \\
\texttt{\#\# Chapitre I}      & \{1: Titre II\}          & [Titre II] \\
\texttt{\#\#\# Article 165}   & \{1: Titre II, 2: Chap. I\} & [Titre II, Chapitre I] \\
\texttt{\#\#\# Article 166}   & \{1: Titre II, 2: Chap. I\} & [Titre II, Chapitre I] \\
\texttt{\#\# Chapitre II}     & \{1: Titre II\}          & [Titre II] \\
\bottomrule
\end{longtable}

\subsection{\texorpdfstring{\texttt{separer\_contenu}}{separer\_contenu} : trier normatif, avertissements et notes (lignes 135 à 187)}
\begin{shaded}
\begin{Verbatim}[numbers=left,firstnumber=135,numbersep=6pt,fontsize=\footnotesize,xleftmargin=1.2em]


# ============================================================================
# SEPARATION NORMATIF / NOTES
# ============================================================================

def separer_contenu(lignes: list[str]) -> tuple[str, list[str], list[str]]:
    """
    Trie le contenu en trois categories :

    1. normatif       -> le texte de loi lui-meme
    2. avertissements -> blocs "> **Avertissement de concordance**"
                         A INDEXER : ils empechent une reponse fausse
                         (renvoi perime de la loi 69-33 vers l'article 161).
    3. notes          -> blocs "> **Note de fidélité**", purement editoriaux,
                         exclus de l'index.
    """
    normatif, avertissements, notes = [], [], []
    bloc, type_bloc = [], None
\end{Verbatim}
\end{shaded}

\begin{explication}
\ligne{135--136} Lignes vides.
\ligne{137--139} Bannière « SEPARATION NORMATIF / NOTES ».
\ligne{140} Ligne vide.
\ligne{141} Définition : reçoit les lignes d'une section et renvoie un \emph{tuple} de trois
éléments : le texte normatif (\Verb@str@), la liste des avertissements et la liste des notes.
\ligne{142--151} Docstring qui décrit les trois catégories :
  \begin{itemize}
    \item \textbf{normatif} : le texte de loi lui-même ;
    \item \textbf{avertissements} : blocs « Avertissement de concordance ». \emph{À indexer} :
          ils empêchent une réponse fausse (le renvoi périmé de la loi 69-33 vers l'article 161) ;
    \item \textbf{notes} : blocs « Note de fidélité », purement éditoriaux, exclus de l'index.
  \end{itemize}
\ligne{152} Trois listes vides.
\ligne{153} \Verb@bloc@ accumule les lignes du bloc \Verb@>@ en cours ; \Verb@type_bloc@ vaut
\Verb@None@ (aucun bloc), \Verb@"avertissement"@ ou \Verb@"note"@.
\end{explication}

\begin{shaded}
\begin{Verbatim}[numbers=left,firstnumber=154,numbersep=6pt,fontsize=\footnotesize,xleftmargin=1.2em]

    def cloturer():
        nonlocal bloc, type_bloc
        if bloc:
            contenu = "\n".join(bloc).strip()
            (avertissements if type_bloc == "avertissement" else notes).append(contenu)
        bloc, type_bloc = [], None
\end{Verbatim}
\end{shaded}

\begin{explication}
\ligne{154} Ligne vide.
\ligne{155} \Verb@cloturer@ est une \emph{fonction imbriquée} (définie dans une autre fonction).
Elle a accès aux variables de \Verb@separer_contenu@ : c'est une \emph{fermeture} (closure).
Elle termine le bloc \Verb@>@ en cours.
\ligne{156} \Verb@nonlocal bloc, type_bloc@ est indispensable, car la ligne 160
\emph{réaffecte} ces variables. Sans \Verb@nonlocal@, Python les considérerait comme locales à
\Verb@cloturer@, et la ligne 157 lèverait \Verb@UnboundLocalError@. Les listes
\Verb@avertissements@ et \Verb@notes@ n'en ont pas besoin : elles sont seulement modifiées
(\Verb@.append@), pas réaffectées.
\ligne{157} Un bloc est-il ouvert ? Une liste non vide est « vraie ».
\ligne{158} On recolle les lignes du bloc avec des retours à la ligne, puis on retire les
espaces aux extrémités.
\ligne{159} L'expression conditionnelle choisit \emph{la liste de destination}, puis on appelle
\Verb@.append@ dessus : avertissement si le type est \Verb@"avertissement"@, note sinon.
\ligne{160} Réinitialisation : plus de bloc en cours.
\end{explication}

\begin{shaded}
\begin{Verbatim}[numbers=left,firstnumber=161,numbersep=6pt,fontsize=\footnotesize,xleftmargin=1.2em]

    for ligne in lignes:
        nue = ligne.strip()

        if nue.startswith(">"):
            contenu = nue.lstrip("> ").rstrip()
            if type_bloc is None:
                type_bloc = ("avertissement" if RE_AVERTISSEMENT.search(contenu)
                             else "note")
            bloc.append(contenu)
            continue
\end{Verbatim}
\end{shaded}

\begin{explication}
\ligne{161} Ligne vide.
\ligne{162} Parcours de chaque ligne de la section.
\ligne{163} \Verb@nue@ est la ligne sans espaces ni \Verb@\r@ aux extrémités. L'indentation
d'origine est perdue.
\ligne{164} Ligne vide.
\ligne{165} Une ligne qui commence par \Verb@>@ est une citation Markdown : elle appartient à un bloc.
\ligne{166} \Verb@nue.lstrip("> ")@ : attention, \Verb@lstrip@ reçoit un \textbf{ensemble de
caractères}, pas un préfixe. Il retire au début tous les \Verb@>@ et espaces, dans
n'importe quel ordre (\Verb@> > texte@ donne \Verb@texte@). \Verb@.rstrip()@ nettoie la fin.
Une ligne composée seulement de \Verb@>@ devient une chaîne vide.
\ligne{167} \Verb@type_bloc is None@ : c'est la \emph{première} ligne du bloc.
\ligne{168--169} Le type du bloc est fixé une fois pour toutes par sa première ligne :
\Verb@RE_AVERTISSEMENT.search@ cherche le motif \emph{n'importe où} dans la ligne. S'il est
trouvé, c'est un avertissement, sinon une note.
\ligne{170} La ligne est ajoutée au bloc.
\ligne{171} \Verb@continue@ : on passe directement à la ligne suivante, sans exécuter la suite de la boucle.
\end{explication}

\begin{shaded}
\begin{Verbatim}[numbers=left,firstnumber=172,numbersep=6pt,fontsize=\footnotesize,xleftmargin=1.2em]

        if bloc:
            cloturer()

        if nue == "---":
            continue
        if nue.startswith("**Source officielle"):
            continue
        if nue.startswith("*(La suite de l'article"):
            continue

        normatif.append(nue)

    cloturer()
    texte = re.sub(r"\n{3,}", "\n\n", "\n".join(normatif)).strip()
    return texte, avertissements, notes
\end{Verbatim}
\end{shaded}

\begin{explication}
\ligne{172} Ligne vide.
\ligne{173--174} La ligne n'est pas une citation mais un bloc était ouvert : le bloc est terminé,
on le clôture. Conséquence : une ligne vide sans \Verb@>@ entre deux paragraphes d'une citation
coupe celle-ci en deux blocs, et le second prend le type de sa propre première ligne.
\ligne{175} Ligne vide.
\ligne{176--177} \Verb@---@ (ligne horizontale Markdown) : ignorée.
\ligne{178--179} La ligne « **Source officielle** » (métadonnée éditoriale) : ignorée.
\ligne{180--181} Marqueur « *(La suite de l'article…* » : ignoré.
\ligne{182} Ligne vide.
\ligne{183} Tout le reste est du texte normatif. Les lignes vides (\Verb@""@) sont aussi
ajoutées : ce sont elles qui gardent la séparation entre paragraphes, utilisée plus tard par
\Verb@sous_decouper@.
\ligne{184} Ligne vide.
\ligne{185} Dernier appel à \Verb@cloturer@ : si la section se termine par un bloc \Verb@>@,
aucune ligne normale n'a déclenché sa clôture.
\ligne{186} On recolle les lignes normatives, puis \Verb@re.sub(r"\n{3,}", "\n\n", ...)@ remplace
toute suite de trois retours à la ligne ou plus par deux : au plus une ligne vide entre
paragraphes. \Verb@.strip()@ retire les lignes vides au début et à la fin.
\ligne{187} Renvoi du tuple (texte, avertissements, notes).
\end{explication}

\subsection{Tableaux et articles trop longs (lignes 188 à 269)}
\begin{shaded}
\begin{Verbatim}[numbers=left,firstnumber=188,numbersep=6pt,fontsize=\footnotesize,xleftmargin=1.2em]


# ============================================================================
# SOUS-DECOUPAGE DES ARTICLES TROP LONGS
# ============================================================================

def eclater_tableau(texte: str) -> list[tuple[str, str]] | None:
    """
    Transforme un tableau markdown en une liste de (etiquette, phrase).

    POURQUOI : le tableau de l'article premier fait 2462 caracteres et
    melange 10 especes. Le signal "sanglier" s'y noyait parmi lievre,
    becasse et tourterelle : un vecteur unique ne peut pas representer
    10 regimes de dates differents.

    Meme principe que l'article 12 : un tableau est une base de donnees
    deguisee en prose, pas de la prose. On le decoupe par LIGNE.

    Retourne None si le texte n'est pas un tableau exploitable.
    """
    lignes = [l for l in texte.splitlines() if RE_LIGNE_TABLEAU.match(l)]
    if len(lignes) < 3:
        return None
\end{Verbatim}
\end{shaded}

\subsubsection*{\texttt{eclater\_tableau} : une ligne de tableau = un chunk}

\begin{explication}
\ligne{188--189} Lignes vides.
\ligne{190--192} Bannière « SOUS-DECOUPAGE DES ARTICLES TROP LONGS ».
\ligne{193} Ligne vide.
\ligne{194} Définition : renvoie une liste de couples \Verb@(etiquette, phrase)@, ou \Verb@None@.
\ligne{195--207} Docstring. Le tableau de l'article premier fait 2462 caractères et mélange 10
espèces. Or un embedding est \textbf{un seul vecteur de taille fixe} : il ne peut pas
représenter dix régimes de dates différents, et le signal « sanglier » s'y noie. Un tableau
est en réalité une base de données écrite sous forme de texte : on le découpe \emph{par ligne}.
\ligne{208} Compréhension de liste : on ne garde que les lignes qui commencent par \Verb@|@.
\ligne{209--210} Moins de 3 lignes (il faut au minimum entête, séparateur et une ligne de
données) : ce n'est pas un tableau exploitable.
\end{explication}

\begin{shaded}
\begin{Verbatim}[numbers=left,firstnumber=211,numbersep=6pt,fontsize=\footnotesize,xleftmargin=1.2em]

    entetes, donnees = None, []
    for l in lignes:
        if RE_LIGNE_SEPARATRICE.match(l):
            continue
        cells = [c.strip() for c in l.strip().strip("|").split("|")]
        if entetes is None:
            entetes = cells
        else:
            donnees.append(cells)

    if not entetes or not donnees:
        return None
\end{Verbatim}
\end{shaded}

\begin{explication}
\ligne{211} Ligne vide.
\ligne{212} \Verb@entetes@ vaut \Verb@None@ tant que la ligne d'entête n'est pas lue ;
\Verb@donnees@ recevra les lignes suivantes.
\ligne{213} Parcours des lignes du tableau.
\ligne{214--215} La ligne \Verb@|---|---|@ est ignorée.
\ligne{216} Découpage d'une ligne en cellules, en trois étapes :
\Verb@l.strip()@ retire les espaces, \Verb@.strip("|")@ retire les barres du début et de la
fin, \Verb@.split("|")@ coupe aux barres restantes ; chaque cellule est ensuite nettoyée par
\Verb@c.strip()@.\\
Exemple : \newline{\footnotesize\Verb@| Sanglier | 5 octobre 2025 | 1er février 2026 |@}\newline  donne
\newline{\footnotesize\Verb@["Sanglier", "5 octobre 2025", "1er février 2026"]@}\newline .
\ligne{217--218} La première ligne non séparatrice est l'entête.
\ligne{219--220} Les suivantes sont des lignes de données.
\ligne{221} Ligne vide.
\ligne{222--223} Sans entête ou sans données : \Verb@None@.
\end{explication}

\begin{shaded}
\begin{Verbatim}[numbers=left,firstnumber=224,numbersep=6pt,fontsize=\footnotesize,xleftmargin=1.2em]

    sortie = []
    for cells in donnees:
        if len(cells) != len(entetes):
            continue
        # Premiere colonne = sujet de la ligne (l'espece)
        sujet = cells[0]
        details = ", ".join(f"{e} : {c}" for e, c in zip(entetes[1:], cells[1:]))
        sortie.append((sujet, f"{sujet}. {details}."))
    return sortie or None
\end{Verbatim}
\end{shaded}

\begin{explication}
\ligne{224} Ligne vide.
\ligne{225} Liste résultat.
\ligne{226} Parcours des lignes de données.
\ligne{227--228} Si le nombre de cellules diffère de celui de l'entête (ligne mal formée), la
ligne est \textbf{ignorée sans message}.
\ligne{229} Commentaire : la première colonne est le sujet de la ligne.
\ligne{230} \Verb@sujet@ = première cellule (l'espèce).
\ligne{231} \Verb@zip(entetes[1:], cells[1:])@ associe chaque entête (sauf la première) à la
cellule correspondante. Pour chaque couple, on écrit \Verb@"entête : valeur"@, puis on joint
le tout avec des virgules.\\
Résultat : \newline{\footnotesize\Verb@Date d'ouverture : 5 octobre 2025, Date de fermeture : 1er février 2026@}\newline .
\ligne{232} On ajoute le couple \Verb@(sujet, phrase)@ ; la phrase vaut \Verb@"sujet. détails."@.
\attention si le sujet finit déjà par un point, on obtient deux points : « (El khedra).. Date
d'ouverture », visible dans \fichier{chunks.json}.\par
\ligne{233} \Verb@sortie or None@ : une liste vide est « fausse », donc on renvoie \Verb@None@
si aucune ligne n'a été retenue.
\end{explication}

\subsubsection*{\texttt{contient\_tableau}}
\begin{shaded}
\begin{Verbatim}[numbers=left,firstnumber=234,numbersep=6pt,fontsize=\footnotesize,xleftmargin=1.2em]


def contient_tableau(texte: str) -> bool:
    """Un tableau markdown ne doit JAMAIS etre coupe au milieu."""
    return any(RE_LIGNE_TABLEAU.match(l) for l in texte.splitlines())
\end{Verbatim}
\end{shaded}

\begin{explication}
\ligne{234--235} Lignes vides.
\ligne{236} Définition : renvoie un booléen.
\ligne{237} Docstring : un tableau Markdown ne doit jamais être coupé au milieu.
\ligne{238} \Verb@any(...)@ renvoie \Verb@True@ dès qu'un élément est vrai, et s'arrête là
(évaluation \emph{court-circuit}). Une seule ligne commençant par \Verb@|@ suffit.
\end{explication}

\subsubsection*{\texttt{sous\_decouper} : couper un article long par paragraphes}
\begin{shaded}
\begin{Verbatim}[numbers=left,firstnumber=239,numbersep=6pt,fontsize=\footnotesize,xleftmargin=1.2em]


def sous_decouper(texte: str) -> list[str]:
    """
    Regroupe les paragraphes en blocs de <= TAILLE_MAX_CHUNK caracteres.

    Pourquoi : un embedding est un vecteur de taille FIXE. Plus le texte est
    long, plus le sens s'y dilue. L'article 3 de l'arrete melange redevances,
    eperviers, faucons, slougui, sanglier et etourneaux : un seul vecteur ne
    peut pas representer tout ca.

    On ne coupe jamais au milieu d'un paragraphe. Et on laisse les tableaux
    entiers meme s'ils depassent : le tableau de l'article premier fait
    ~2500 caracteres mais forme un bloc coherent.
    """
    if contient_tableau(texte) or len(texte) <= config.TAILLE_MAX_CHUNK:
        return [texte]

    paragraphes = [p.strip() for p in re.split(r"\n\s*\n", texte) if p.strip()]
    blocs, courant = [], []
\end{Verbatim}
\end{shaded}

\begin{explication}
\ligne{239--240} Lignes vides.
\ligne{241} Définition : renvoie une liste de blocs de texte.
\ligne{242--253} Docstring. Un embedding est un vecteur de \textbf{taille fixe} (1024 nombres) :
plus le texte est long, plus son sens est dilué. L'article 3 de l'arrêté mélange redevances,
éperviers, faucons, slougui, sanglier et étourneaux : un seul vecteur ne peut pas bien
représenter tout cela. On ne coupe jamais au milieu d'un paragraphe, et on laisse les tableaux
entiers.
\ligne{254} Si le texte contient un tableau, ou s'il est assez court
($\leq$ \Verb@TAILLE_MAX_CHUNK@ = 1500 caractères)…
\ligne{255} … on le renvoie tel quel, dans une liste à un seul élément. Renvoyer toujours une
liste simplifie le code appelant.
\ligne{256} Ligne vide.
\ligne{257} \Verb@re.split(r"\n\s*\n", texte)@ découpe aux lignes vides (qui peuvent contenir des
espaces) : on obtient les paragraphes. La compréhension les nettoie et élimine les vides.
\ligne{258} \Verb@blocs@ = résultat ; \Verb@courant@ = paragraphes du bloc en construction.
\end{explication}

\begin{shaded}
\begin{Verbatim}[numbers=left,firstnumber=259,numbersep=6pt,fontsize=\footnotesize,xleftmargin=1.2em]

    for p in paragraphes:
        if courant and sum(len(x) for x in courant) + len(p) > config.TAILLE_MAX_CHUNK:
            blocs.append("\n\n".join(courant))
            courant = [p]
        else:
            courant.append(p)

    if courant:
        blocs.append("\n\n".join(courant))
    return blocs
\end{Verbatim}
\end{shaded}

\begin{explication}
\ligne{259} Ligne vide.
\ligne{260} Parcours des paragraphes.
\ligne{261} Condition de fermeture du bloc courant :
  \begin{itemize}
    \item \Verb@courant@ n'est pas vide : sans cette condition, un paragraphe seul de plus de
          1500 caractères produirait un bloc vide ; avec elle, il est accepté seul dans son bloc ;
    \item \textbf{et} la taille du bloc plus celle du nouveau paragraphe dépasserait 1500.
  \end{itemize}
  La somme ne compte pas les \Verb@"\n\n"@ ajoutés au recollage : un bloc peut dépasser
  1500 de quelques caractères. La somme est recalculée à chaque tour : c'est un coût
  quadratique, négligeable ici.
\ligne{262} Le bloc courant est fermé : paragraphes joints par une ligne vide.
\ligne{263} Le nouveau paragraphe ouvre le bloc suivant.
\ligne{264--265} Sinon, le paragraphe s'ajoute au bloc courant.
\ligne{266} Ligne vide.
\ligne{267--268} Le dernier bloc en construction est ajouté.
\ligne{269} Renvoi des blocs.
\end{explication}

\subsection{Préfixes contextuels (lignes 270 à 327)}
\begin{shaded}
\begin{Verbatim}[numbers=left,firstnumber=270,numbersep=6pt,fontsize=\footnotesize,xleftmargin=1.2em]


# ============================================================================
# PREFIXE CONTEXTUEL
# ============================================================================

def _libelle_court(meta: dict) -> str:
    """
    Version abregee de la source, pour le prefixe d'embedding.
    "Arrêté du 27 août 2025 (JORT n° 107 du 29 août 2025)" -> "Arrêté chasse 2025/2026"
    """
    sid = meta.get("source_id", "")
    return {
        "arrete_2025": "Arrêté chasse 2025/2026",
        "code_forestier": "Code forestier",
        "loi_69_33": "Loi armes 69-33",
        "code_forestier_camper": "Code forestier (camping)",
        "arretes_parcs_1984": "Arrêtés parcs nationaux 1984",
        "production_bio_apiculture": "Apiculture bio (arrêté 2005)",
        "preparation_bio_miel": "Préparation miel bio (arrêté 2005)",
        "guide_technique_apiculture": "Guide technique apiculture",
    }.get(sid, meta["source"][:40])
\end{Verbatim}
\end{shaded}

Un article isolé (« Il est interdit de… ») ne dit pas de quelle loi il vient. On lui ajoute
donc un \textbf{préfixe}. Le code en construit \textbf{deux versions}, pour deux usages différents.

\subsubsection*{\texttt{\_libelle\_court}}

\begin{explication}
\ligne{270--271} Lignes vides.
\ligne{272--274} Bannière « PREFIXE CONTEXTUEL ».
\ligne{275} Ligne vide.
\ligne{276} Le \Verb@_@ initial est une convention Python : fonction \emph{interne} au module,
à ne pas importer ailleurs.
\ligne{277--280} Docstring : la source longue (« Arrêté du 27 août 2025 (JORT n° 107…) ») devient
« Arrêté chasse 2025/2026 ».
\ligne{281} \Verb@meta.get("source_id", "")@ : lecture sans risque de \Verb@KeyError@ ; renvoie
\Verb@""@ si la clé est absente.
\ligne{282} Début d'un dictionnaire littéral, créé à chaque appel.
\ligne{283--290} Table de correspondance \Verb@source_id@ $\rightarrow$ libellé court, une entrée
par document du corpus.
\ligne{291} \Verb@.get(sid, meta["source"][:40])@ : on cherche le libellé ; s'il est absent, on
prend les 40 premiers caractères de la source. Pour tout nouveau document, il faut penser à
ajouter une entrée ici, sinon le libellé est tronqué.
\end{explication}

\subsubsection*{\texttt{construire\_prefixe} : le préfixe LONG}
\begin{shaded}
\begin{Verbatim}[numbers=left,firstnumber=292,numbersep=6pt,fontsize=\footnotesize,xleftmargin=1.2em]


def construire_prefixe(meta: dict, parents: list[str], titre: str) -> str:
    """
    PREFIXE LONG — pour BM25 uniquement, et pour l'affichage.

    Il porte la source complete et toute la hierarchie. Sur un index lexical
    c'est gratuit : BM25 pondere par l'IDF, donc les mots repetes dans les
    160 chunks ("arrêté", "titre") ont un poids quasi nul automatiquement.
    """
    morceaux = [meta["source"]]
    for p in parents:
        if p.startswith(("Titre", "Chapitre", "Partie", "Section")):
            morceaux.append(p)
    morceaux.append(titre)
    return " — ".join(morceaux)
\end{Verbatim}
\end{shaded}

\begin{explication}
\ligne{292--293} Lignes vides.
\ligne{294} Définition : reçoit les métadonnées, la pile des parents et le titre.
\ligne{295--301} Docstring. Le préfixe long sert à BM25 et à l'affichage. Sur un index lexical,
il ne coûte rien : BM25 pondère chaque mot par son \textbf{IDF} (\emph{Inverse Document
Frequency}). Un mot présent dans presque tous les chunks (« arrêté », « titre ») a un poids
quasi nul.
\ligne{302} La liste commence par la source complète.
\ligne{303} Parcours des parents, du plus général au plus précis.
\ligne{304} \Verb@str.startswith@ accepte un \emph{tuple} de préfixes : vrai si l'un d'eux
correspond. On ne garde que les parents \emph{structurels} (Titre, Chapitre, Partie,
Section). Sont exclus le titre H1 du document et un parent « Article 12 ».
\ligne{305} Ajout du parent retenu.
\ligne{306} Ajout du titre de la section.
\ligne{307} Assemblage avec « — ».\\
Exemple : « Arrêté du 27 août 2025 (JORT…) — Titre premier : Réglementation générale —
Article premier ».
\end{explication}

\subsubsection*{\texttt{construire\_prefixe\_court} : le préfixe COURT}
\begin{shaded}
\begin{Verbatim}[numbers=left,firstnumber=308,numbersep=6pt,fontsize=\footnotesize,xleftmargin=1.2em]


def construire_prefixe_court(meta: dict, titre: str, extra: str = "") -> str:
    """
    PREFIXE COURT — pour l'EMBEDDING.

    POURQUOI DEUX PREFIXES :
    un embedding est un vecteur de taille fixe ou CHAQUE token compte. Le
    prefixe long fait ~25 tokens IDENTIQUES sur les 57 chunks de l'arrete.
    Sur un chunk court, il representait la moitie du contenu : tous les
    vecteurs se ressemblaient et la partie discriminante pesait trop peu.
    C'est ce qui faisait remonter l'article 10 (colportage) en tete de trois
    questions sans rapport.

    Contrairement a BM25, l'embedding n'a AUCUN mecanisme equivalent a l'IDF
    pour ignorer ce qui est commun a tout le corpus. Il faut donc le lui
    epargner : ici ~6 tokens au lieu de 25.
    """
    base = f"{_libelle_court(meta)}, {titre}"
    return f"{base}. {extra}" if extra else base
\end{Verbatim}
\end{shaded}

\begin{explication}
\ligne{308--309} Lignes vides.
\ligne{310} Définition ; \Verb@extra@ est facultatif (chaîne vide par défaut).
\ligne{311--325} Docstring, la justification la plus importante du fichier. Dans un embedding,
\textbf{chaque token compte}. Le préfixe long ajoutait environ 25 tokens identiques aux 57
chunks de l'arrêté. Sur un chunk court, cela représentait la moitié du contenu : tous les
vecteurs se ressemblaient, et c'est ce qui faisait remonter l'article 10 (colportage) pour
trois questions sans rapport. Contrairement à BM25, l'embedding \textbf{n'a aucun mécanisme
équivalent à l'IDF} pour ignorer ce qui est commun à tout le corpus : il faut donc le lui
épargner (environ 6 tokens au lieu de 25).
\ligne{326} \Verb@base@ = « libellé court, titre ».
\ligne{327} Si \Verb@extra@ n'est pas vide, on l'ajoute après un point (le nom du gouvernorat,
l'espèce d'une ligne de tableau…).
\end{explication}

\begin{longtable}{p{3cm} p{6cm} p{6cm}}
\toprule
 & \textbf{Préfixe long} & \textbf{Préfixe court} \\
\midrule
Utilisé par & BM25 (\texttt{texte\_indexe}) et affichage & embedding (\texttt{texte\_embedding}) \\
Contenu & source complète + Titre/Chapitre + titre & libellé court + titre (+ extra) \\
Pourquoi & l'IDF neutralise les mots répétés & aucun IDF : les mots répétés diluent le vecteur \\
\bottomrule
\end{longtable}

\subsection{\texorpdfstring{\texttt{charger\_manifeste}}{charger\_manifeste} (lignes 328 à 345)}
\begin{shaded}
\begin{Verbatim}[numbers=left,firstnumber=328,numbersep=6pt,fontsize=\footnotesize,xleftmargin=1.2em]


# ============================================================================
# CHARGEMENT DES METADONNEES D'UN DOMAINE
# ============================================================================

def charger_manifeste(dossier_domaine: Path) -> dict:
    """
    Lit data/corpus/<domaine>/_manifest.json s'il existe.

    Le nom du dossier donne le DOMAINE, mais pas le rang normatif ni la
    validite : c'est le role du manifeste. S'il est absent (cas de camping/
    ou app/), on applique META_DEFAUT et tout fonctionne quand meme.
    """
    f = dossier_domaine / "_manifest.json"
    if f.exists():
        return json.loads(f.read_text(encoding="utf-8"))
    return {}
\end{Verbatim}
\end{shaded}

\begin{explication}
\ligne{328--329} Lignes vides.
\ligne{330--332} Bannière « CHARGEMENT DES METADONNEES D'UN DOMAINE ».
\ligne{333} Ligne vide.
\ligne{334} Définition : reçoit le dossier d'un domaine (\Verb@Path@), renvoie un dictionnaire.
\ligne{335--341} Docstring. Le nom du dossier donne le domaine, mais pas le rang normatif ni la
validité : c'est le rôle du manifeste. S'il est absent, \Verb@META_DEFAUT@ s'applique.
\ligne{342} L'opérateur \Verb@/@ est redéfini par \Verb@Path@ pour \emph{joindre des chemins} :
\Verb@data/corpus/chasse@ \Verb@/@ \Verb@"_manifest.json"@.
\ligne{343} Le fichier existe-t-il ?
\ligne{344} Lecture en UTF-8, puis \Verb@json.loads@ convertit le texte JSON en dictionnaire
Python. Le manifeste est indexé par \emph{nom de fichier} :
\newline{\footnotesize\Verb@{"code_forestier_chasse.md": {"source": ..., "rang_normatif": "loi", ...}}@}\newline .
Un JSON mal formé lève une exception et interrompt tout le chunking (pas de \Verb@try@).
\ligne{345} Pas de manifeste : dictionnaire vide.
\end{explication}

\subsection{\texorpdfstring{\texttt{chunker\_fichier}}{chunker\_fichier} : le cœur du découpage (lignes 346 à 438)}
\begin{shaded}
\begin{Verbatim}[numbers=left,firstnumber=346,numbersep=6pt,fontsize=\footnotesize,xleftmargin=1.2em]


# ============================================================================
# TRAITEMENT D'UN FICHIER
# ============================================================================

def chunker_fichier(chemin: Path, domaine: str, meta_fichier: dict) -> list[dict]:
    texte = chemin.read_text(encoding="utf-8")
    chunks = []

    for sec in decouper_sections(texte):
        corps, avertissements, _ = separer_contenu(sec["lignes"])
        titre = sec["titre"]
        prefixe = construire_prefixe(meta_fichier, sec["parents"], titre)
\end{Verbatim}
\end{shaded}

Pour chaque section, \textbf{un seul cas} s'applique : gouvernorat, section non-article, ou
article (avec ou sans tableau). Les avertissements sont traités en plus, pour les articles.

\begin{shaded}
\begin{Verbatim}[fontsize=\footnotesize]
section
  +-- EXCEPTION 1 : titre "Gouvernorat de X"   -> 1 chunk "zone_geographique"
  +-- EXCEPTION 2 : pas un article (Visas...)  -> 1 chunk "contexte" si > 200 caracteres
  +-- CAS NORMAL  : un article
        +-- EXCEPTION 4 : contient un tableau  -> 1 chunk "tableau" par ligne
        |                                         + 1 chunk "norme" (texte hors tableau)
        +-- sinon                              -> 1 ou N chunks "norme" (sous_decouper)
        +-- EXCEPTION 3 (en plus) : chaque avertissement -> 1 chunk "avertissement"
\end{Verbatim}
\end{shaded}

\begin{explication}
\ligne{346--347} Lignes vides.
\ligne{348--350} Bannière « TRAITEMENT D'UN FICHIER ».
\ligne{351} Ligne vide.
\ligne{352} Définition : chemin du fichier, nom du domaine, métadonnées du fichier. Pas de docstring.
\ligne{353} Lecture du fichier entier en UTF-8.
\ligne{354} Liste des chunks produits pour ce fichier.
\ligne{355} Ligne vide.
\ligne{356} Boucle sur les sections produites par \Verb@decouper_sections@.
\ligne{357} Déballage du tuple renvoyé par \Verb@separer_contenu@. Le troisième élément (les
notes) est affecté à \Verb@_@ : convention Python pour « valeur volontairement ignorée ».
Les notes de fidélité disparaissent donc ici.
\ligne{358} Raccourci vers le titre de la section.
\ligne{359} Préfixe long calculé pour toutes les sections, même celles qui seront ignorées.
\end{explication}

\begin{shaded}
\begin{Verbatim}[numbers=left,firstnumber=360,numbersep=6pt,fontsize=\footnotesize,xleftmargin=1.2em]

        # ---- EXCEPTION 1 : sous-sections geographiques de l'article 12 ----
        m_gouv = RE_GOUVERNORAT.match(titre)
        if m_gouv and corps:
            # "### Gouvernorat de Siliana" n'est pas un article : c'est une
            # subdivision de l'article 12. On remonte la pile des parents
            # pour retrouver l'article de rattachement.
            art_parent = next((numero_article(p) for p in reversed(sec["parents"])
                               if numero_article(p)), None)
            gouv = m_gouv.group(1).strip()
            court = construire_prefixe_court(
                meta_fichier, f"Article {art_parent}",
                f"Réserves de chasse interdites, gouvernorat de {gouv}")
            chunks.append(_fabriquer(meta_fichier, domaine, art_parent,
                                     slug(gouv), titre, prefixe, court, corps,
                                     "zone_geographique", gouv))
            continue
\end{Verbatim}
\end{shaded}

\begin{explication}
\ligne{360} Ligne vide.
\ligne{361} Commentaire : exception 1, les sous-sections géographiques de l'article 12.
\ligne{362} Le titre est-il « Gouvernorat de X » ?
\ligne{363} Oui, et la section a du contenu : on la traite.
\ligne{364--366} Commentaires : « \#\#\# Gouvernorat de Siliana » n'est pas un article mais une
subdivision de l'article 12, qu'il faut retrouver dans la pile des parents.
\ligne{367--368} \Verb@next(générateur, None)@ renvoie le premier élément produit par le
générateur, ou \Verb@None@ s'il n'en produit aucun. Le générateur parcourt les parents
\textbf{du plus proche au plus lointain} (\Verb@reversed@) et ne produit que ceux qui sont des
articles. Résultat : \Verb@art_parent = "12"@. (\Verb@numero_article@ est appelé deux fois
par parent : petite redondance.)
\ligne{369} Nom du gouvernorat : groupe 1 de la regex, sans espaces.
\ligne{370--372} Préfixe court avec titre « Article 12 » et extra « Réserves de chasse interdites,
gouvernorat de Siliana ». Ce texte est écrit en dur : toute section « Gouvernorat de… » de
n'importe quel domaine recevrait ce libellé.
\ligne{373--375} Création du chunk par \Verb@_fabriquer@. Arguments, dans l'ordre : métadonnées,
domaine, article (\Verb@"12"@), sous-clé \Verb@slug(gouv)@, titre, préfixe long, préfixe
court, texte (\Verb@corps@), type \Verb@"zone_geographique"@, gouvernorat.
Identifiant obtenu : \Verb@arrete_2025_art12_siliana@ (24 chunks de ce type dans le corpus).
\ligne{376} \Verb@continue@ : section traitée, on passe à la suivante.
\end{explication}

\begin{shaded}
\begin{Verbatim}[numbers=left,firstnumber=377,numbersep=6pt,fontsize=\footnotesize,xleftmargin=1.2em]

        article = numero_article(titre)

        # ---- EXCEPTION 2 : sections qui ne sont pas des articles ----------
        if article is None:
            # Visas, Preambule, Lacunes connues. Non normatifs mais utiles.
            # "Lacunes connues" permet au bot de repondre "cette information
            # n'est pas dans mon corpus" au lieu d'inventer.
            if corps and len(corps) > config.TAILLE_MIN_ALERTE:
                court = construire_prefixe_court(meta_fichier, titre)
                chunks.append(_fabriquer(meta_fichier, domaine, None,
                                         slug(titre), titre, prefixe, court,
                                         corps, "contexte", None))
            continue
\end{Verbatim}
\end{shaded}

\begin{explication}
\ligne{377} Ligne vide.
\ligne{378} Numéro d'article du titre, ou \Verb@None@.
\ligne{379} Ligne vide.
\ligne{380} Commentaire : exception 2, sections qui ne sont pas des articles.
\ligne{381} Ce n'est pas un article…
\ligne{382--384} Commentaires : Visas, Préambule, « Lacunes connues ». Ces sections ne sont pas
normatives mais utiles : « Lacunes connues » permet au bot de répondre « cette information
n'est pas dans mon corpus » au lieu d'inventer.
\ligne{385} … on ne crée un chunk que si le texte dépasse \Verb@TAILLE_MIN_ALERTE@ (200 caractères).
Cela élimine les titres vides (« \#\# Chapitre II » sans texte propre).
\ligne{386} Préfixe court avec le titre de la section.
\ligne{387--389} Chunk de type \Verb@"contexte"@, sans article. La sous-clé est
\Verb@slug(titre)@. Identifiant réel : \Verb@arrete_2025_artvisas_visas@.
\ligne{390} \Verb@continue@ : une section trop courte est ignorée \textbf{sans message}.
\end{explication}
\attention c'est aussi le cas des documents sans articles (\fichier{app/}, guides
d'apiculture) : chacune de leurs sections est un « contexte ». Vérifié sur
\fichier{aide\_application.md} : trois réponses de la FAQ font moins de 200 caractères et
\textbf{ne sont jamais indexées} : « Puis-je créer une alerte sans réseau ? » (148),
« Comment suivre mes signalements ? » (173), « Les réponses du chatbot ont-elles une valeur
légale ? » (160). Le chatbot ne peut donc pas répondre à ces questions.\par

\begin{shaded}
\begin{Verbatim}[numbers=left,firstnumber=391,numbersep=6pt,fontsize=\footnotesize,xleftmargin=1.2em]

        # ---- CAS NORMAL : un article --------------------------------------
        if corps:
            # --- EXCEPTION 4 : article contenant un tableau ----------------
            # Une ligne de tableau = un chunk. Sinon les 10 especes du
            # tableau de l'article premier se noient dans un seul vecteur.
            lignes_tab = eclater_tableau(corps) if contient_tableau(corps) else None
\end{Verbatim}
\end{shaded}

\begin{explication}
\ligne{391} Ligne vide.
\ligne{392} Commentaire : cas normal, un article.
\ligne{393} L'article a-t-il du texte ? Un article sans corps ne produit pas de chunk
« norme », mais ses avertissements sont traités quand même (ligne 429, hors de ce \Verb@if@).
\ligne{394--396} Commentaires : exception 4, une ligne de tableau = un chunk.
\ligne{397} Expression conditionnelle : on n'appelle \Verb@eclater_tableau@ que si le texte
contient un tableau ; sinon \Verb@lignes_tab = None@.
\end{explication}

\begin{shaded}
\begin{Verbatim}[numbers=left,firstnumber=398,numbersep=6pt,fontsize=\footnotesize,xleftmargin=1.2em]

            if lignes_tab:
                # le texte hors tableau (alineas qui suivent) devient un chunk
                hors_tableau = "\n".join(
                    l for l in corps.splitlines()
                    if not RE_LIGNE_TABLEAU.match(l)).strip()

                for sujet, phrase in lignes_tab:
                    court = construire_prefixe_court(meta_fichier, titre, sujet)
                    chunks.append(_fabriquer(
                        meta_fichier, domaine, article, slug(sujet, 30), titre,
                        f"{prefixe} — {sujet}", court, phrase, "tableau", None))

                if len(hors_tableau) > config.TAILLE_MIN_ALERTE:
                    court = construire_prefixe_court(meta_fichier, titre)
                    chunks.append(_fabriquer(
                        meta_fichier, domaine, article, "suite", titre,
                        prefixe, court, hors_tableau, "norme", None))
\end{Verbatim}
\end{shaded}

\begin{explication}
\ligne{398} Ligne vide.
\ligne{399} Le tableau est exploitable (liste non vide).
\ligne{400} Commentaire : le texte hors tableau devient un chunk.
\ligne{401--403} \Verb@hors_tableau@ : on garde les lignes du corps qui \emph{ne} sont \emph{pas}
des lignes de tableau (expression génératrice passée à \Verb@join@), puis on nettoie.
\ligne{404} Ligne vide.
\ligne{405} Pour chaque ligne du tableau : déballage du couple \Verb@(sujet, phrase)@.
\ligne{406} Préfixe court avec \Verb@extra = sujet@ : l'espèce apparaît dans le préfixe
\emph{et} dans la phrase, ce qui renforce son poids dans le vecteur.
\ligne{407--409} Chunk de type \Verb@"tableau"@ : sous-clé \Verb@slug(sujet, 30)@ (tronquée à
30 caractères), préfixe long complété par « — sujet », texte = la phrase.
Identifiant réel : \newline{\footnotesize\Verb@arrete_2025_art1_lievre_perdrix_caille_sedentai@}\newline . Deux lignes dont
les 30 premiers caractères seraient identiques produiraient le même identifiant.
\ligne{410} Ligne vide.
\ligne{411} Le texte hors tableau fait plus de 200 caractères…
\ligne{412} … préfixe court sans extra…
\ligne{413--415} … et chunk \Verb@"norme"@ de sous-clé \Verb@"suite"@
(\Verb@arrete_2025_art1_suite@). En dessous de 200 caractères, ce texte est perdu.
\end{explication}

\begin{shaded}
\begin{Verbatim}[numbers=left,firstnumber=416,numbersep=6pt,fontsize=\footnotesize,xleftmargin=1.2em]
            else:
                blocs = sous_decouper(corps)
                for i, bloc in enumerate(blocs, start=1):
                    # Si l'article a ete sous-decoupe, on l'indique au LLM
                    # pour qu'il sache qu'il ne voit qu'un extrait.
                    suffixe = f" (partie {i}/{len(blocs)})" if len(blocs) > 1 else ""
                    court = construire_prefixe_court(meta_fichier, titre)
                    chunks.append(_fabriquer(meta_fichier, domaine, article,
                                             f"p{i}" if len(blocs) > 1 else None,
                                             titre, prefixe + suffixe, court,
                                             bloc, "norme", None))
\end{Verbatim}
\end{shaded}

\begin{explication}
\ligne{416} Pas de tableau, ou tableau mal formé (\Verb@eclater_tableau@ a renvoyé \Verb@None@).
Dans ce dernier cas, \Verb@sous_decouper@ renvoie le texte entier sans le couper.
\ligne{417} Découpage éventuel en blocs de 1500 caractères au maximum.
\ligne{418} \Verb@enumerate(blocs, start=1)@ fournit l'indice \Verb@i@ (à partir de 1) et le bloc.
\ligne{419--420} Commentaires : si l'article a été découpé, on l'indique au LLM.
\ligne{421} \Verb@suffixe@ = « (partie 2/3) » s'il y a plusieurs blocs, chaîne vide sinon.
\ligne{422} Préfixe court, identique pour toutes les parties.
\ligne{423--426} Chunk \Verb@"norme"@ : sous-clé \Verb@"p1"@, \Verb@"p2"@… s'il y a plusieurs
blocs, sinon \Verb@None@. Le suffixe est ajouté au préfixe long.
Identifiants réels : \Verb@code_forestier_art165@ (un bloc), \Verb@arrete_2025_art3_p1@ (plusieurs).
\end{explication}

\begin{shaded}
\begin{Verbatim}[numbers=left,firstnumber=427,numbersep=6pt,fontsize=\footnotesize,xleftmargin=1.2em]

        # ---- EXCEPTION 3 : avertissements de concordance -------------------
        for j, avert in enumerate(avertissements, start=1):
            court = construire_prefixe_court(
                meta_fichier, titre, "avertissement de concordance")
            chunks.append(_fabriquer(meta_fichier, domaine, article,
                                     f"avert{j}",
                                     f"{titre} — avertissement de concordance",
                                     prefixe + " — avertissement de concordance",
                                     court, avert, "avertissement", None))

    return chunks
\end{Verbatim}
\end{shaded}

\begin{explication}
\ligne{427} Ligne vide.
\ligne{428} Commentaire : exception 3, les avertissements de concordance.
\ligne{429} Parcours des avertissements, numérotés à partir de 1. Cette boucle est au niveau de
\Verb@if corps@ : elle s'exécute pour tout article, même sans corps.
\ligne{430--431} Préfixe court avec l'extra « avertissement de concordance ».
\ligne{432--436} Chunk \Verb@"avertissement"@ : sous-clé \Verb@avert1@, \Verb@avert2@… ; le titre
et le préfixe long sont complétés par « — avertissement de concordance ».
Identifiants réels : \Verb@code_forestier_art161_avert1@ et \Verb@loi_69_33_art16_avert1@.
Le type \Verb@"avertissement"@ fera ajouter l'étiquette \Verb@[AVERTISSEMENT DE CONCORDANCE]@
dans le prompt (\fichier{generator.py}).
\ligne{437} Ligne vide.
\ligne{438} Renvoi des chunks du fichier.
\end{explication}
\attention un avertissement placé dans une section « Gouvernorat » ou non-article est perdu,
car les \Verb@continue@ des lignes 376 et 390 sautent cette boucle. Aucun cas actuel dans le corpus.\par

\subsection{\texorpdfstring{\texttt{\_fabriquer}}{\_fabriquer} : la forme d'un chunk (lignes 439 à 479)}
\begin{shaded}
\begin{Verbatim}[numbers=left,firstnumber=439,numbersep=6pt,fontsize=\footnotesize,xleftmargin=1.2em]


def _fabriquer(meta_fichier, domaine, article, sous_cle, titre, prefixe,
               prefixe_court, texte, type_chunk, gouvernorat) -> dict:
    """
    Fabrique un chunk normalise.

    Les metadonnees restent des types simples (str, int, float, bool) :
    ni liste, ni dict. D'ou les alias joints en une seule chaine. Cela
    garde le format lisible en JSON et compatible avec n'importe quel
    magasin de vecteurs.
    """
    ident = f"{meta_fichier['source_id']}_art{article or slug(titre, 20)}"
    if sous_cle:
        ident += f"_{sous_cle}"

\end{Verbatim}
\end{shaded}

\begin{explication}
\ligne{439--440} Lignes vides.
\ligne{441--442} Définition sur deux lignes (continuation implicite dans les parenthèses), avec
10 paramètres positionnels et sans annotations de type.
\ligne{443--450} Docstring. Les métadonnées ne contiennent que des types simples (\Verb@str@,
\Verb@int@, \Verb@float@, \Verb@bool@), jamais de listes ni de dictionnaires. C'est pourquoi
les alias sont joints en une seule chaîne. Le format reste lisible en JSON et compatible
avec n'importe quelle base vectorielle (ChromaDB, par exemple, refuse les listes en métadonnées).
\ligne{451} Identifiant : \Verb@source_id@ + \Verb@_art@ + numéro d'article. Sans article,
\Verb@article or slug(titre, 20)@ prend le slug du titre (20 caractères).
\ligne{452--453} Si une sous-clé existe (\Verb@siliana@, \Verb@p2@, \Verb@suite@, \Verb@avert1@…),
on l'ajoute avec \Verb@+=@ (concaténation). Pour une section non-article, le titre apparaît
deux fois : \newline{\footnotesize\Verb@aide_application_artcomment_modifier_mon_comment_modifier_mon_profil@}\newline .
\ligne{454} Ligne vide.
\end{explication}

\begin{shaded}
\begin{Verbatim}[numbers=left,firstnumber=455,numbersep=6pt,fontsize=\footnotesize,xleftmargin=1.2em]
    return {
        "id": ident,
        # texte_embedding -> COURT, ce qu'on vectorise. Chaque token compte.
        "texte_embedding": f"{prefixe_court}\n{texte}",
        # texte_indexe -> LONG, ce que BM25 lit. L'IDF neutralise le commun.
        "texte_indexe": f"{prefixe}\n\n{texte}",
        # texte_affiche -> ce qu'on MONTRE au LLM et a l'utilisateur
        "texte_affiche": texte,
        "prefixe": prefixe,
        "prefixe_court": prefixe_court,
        "metadata": {
            "domaine": domaine,
            "source": meta_fichier["source"],
            "source_id": meta_fichier["source_id"],
            "rang_normatif": meta_fichier["rang_normatif"],
            "rang_poids": meta_fichier["rang_poids"],
            "validite": meta_fichier["validite"],
            "article": article or "",
            "titre": titre,
            "type": type_chunk,
            "gouvernorat": gouvernorat or "",
            "alias": " | ".join(extraire_alias_sic(texte)),
            "nb_caracteres": len(texte),
        },
    }
\end{Verbatim}
\end{shaded}

\begin{explication}
\ligne{455} Renvoi d'un dictionnaire littéral.
\ligne{456} \Verb@"id"@ : identifiant unique du chunk. Il relie \fichier{chunks.json},
\fichier{vecteurs.npy} (même ordre) et \fichier{bm25.pkl}.
\ligne{457} Commentaire.
\ligne{458} \Verb@"texte_embedding"@ = préfixe \textbf{court} + retour à la ligne + texte. C'est
ce que \fichier{indexer.py} transforme en vecteur.
\ligne{459} Commentaire.
\ligne{460} \Verb@"texte_indexe"@ = préfixe \textbf{long} + ligne vide + texte. C'est ce que BM25 lit.
\ligne{461} Commentaire.
\ligne{462} \Verb@"texte_affiche"@ = le texte seul. C'est ce qu'on envoie au LLM et qu'on montre
à l'utilisateur (via \fichier{store.py}).
\ligne{463} \Verb@"prefixe"@ : préfixe long, conservé pour le débogage.
\ligne{464} \Verb@"prefixe_court"@ : préfixe court, conservé pour le débogage.
\ligne{465} Ouverture du sous-dictionnaire \Verb@"metadata"@.
\ligne{466} \Verb@domaine@ : nom du dossier (\Verb@chasse@, \Verb@camper@…). Sert au
\textbf{filtrage par spécialité} dans \fichier{store.py} et \fichier{retriever.py}.
\ligne{467} \Verb@source@ : nom complet du texte, cité par le LLM et affiché dans les sources.
\ligne{468} \Verb@source_id@ : identifiant court du texte.
\ligne{469} \Verb@rang_normatif@ (\Verb@loi@, \Verb@arrete@, \Verb@documentaire@) : affiché dans
l'en-tête de l'extrait, il permet au LLM d'appliquer la hiérarchie des normes (règle 4 du prompt).
\ligne{470} \Verb@rang_poids@ : version numérique du rang (1 = loi, 2 = arrêté, 9 = documentaire).
Il n'est lu par aucun autre fichier pour l'instant.
\ligne{471} \Verb@validite@ (\Verb@permanent@, \Verb@saison_2025_2026@) : permet au LLM de
signaler une période écoulée (règle 3 du prompt).
\ligne{472} \Verb@article@ : numéro ou \Verb@""@. On utilise une chaîne vide plutôt que
\Verb@None@ pour garder des types simples.
\ligne{473} \Verb@titre@ : titre de la section.
\ligne{474} \Verb@type@ : \Verb@norme@, \Verb@tableau@, \Verb@zone_geographique@,
\Verb@contexte@ ou \Verb@avertissement@. \fichier{main.py} n'affiche comme sources que les
trois premiers.
\ligne{475} \Verb@gouvernorat@ : nom ou \Verb@""@.
\ligne{476} \Verb@alias@ : les formes correctes des coquilles, jointes par \Verb@" | "@.
\fichier{indexer.py} les ajoute au texte indexé par BM25.
\ligne{477} \Verb@nb_caracteres@ : longueur du texte affiché, en \emph{caractères} (pas en tokens).
Sert au rapport de contrôle.
\ligne{478} Fermeture de \Verb@metadata@.
\ligne{479} Fermeture du chunk.
\end{explication}

\subsection{\texorpdfstring{\texttt{chunker\_corpus}}{chunker\_corpus} : le point d'entrée (lignes 480 à 523)}
\begin{shaded}
\begin{Verbatim}[numbers=left,firstnumber=480,numbersep=6pt,fontsize=\footnotesize,xleftmargin=1.2em]


# ============================================================================
# POINT D'ENTREE : SCAN DE TOUS LES DOMAINES
# ============================================================================

def chunker_corpus(verbeux: bool = True) -> list[dict]:
    """
    Parcourt data/corpus/<domaine>/*.md et produit tous les chunks.
    Le nom du sous-dossier devient la metadonnee "domaine".
    """
    tous = []
\end{Verbatim}
\end{shaded}

\begin{explication}
\ligne{480--481} Lignes vides.
\ligne{482--484} Bannière « POINT D'ENTREE : SCAN DE TOUS LES DOMAINES ».
\ligne{485} Ligne vide.
\ligne{486} Définition : \Verb@verbeux@ (affichage de la progression) vaut \Verb@True@ par défaut.
C'est la fonction appelée par \fichier{scripts/build\_embeddings.py}.
\ligne{487--490} Docstring : parcourt \fichier{data/corpus/<domaine>/*.md} ; le nom du
sous-dossier devient la métadonnée \Verb@domaine@.
\ligne{491} \Verb@tous@ accumulera les chunks de tous les fichiers.
\end{explication}

\begin{shaded}
\begin{Verbatim}[numbers=left,firstnumber=492,numbersep=6pt,fontsize=\footnotesize,xleftmargin=1.2em]

    for dossier in sorted(config.CORPUS_DIR.iterdir()):
        if not dossier.is_dir():
            continue
        domaine = dossier.name
        manifeste = charger_manifeste(dossier)
        fichiers = sorted(dossier.glob("*.md"))

        if not fichiers:
            continue
        if verbeux:
            print(f"\n  [{domaine}]")
\end{Verbatim}
\end{shaded}

\begin{explication}
\ligne{492} Ligne vide.
\ligne{493} \Verb@CORPUS_DIR.iterdir()@ liste le contenu de \fichier{data/corpus/}. L'ordre
renvoyé par le système n'est pas garanti : \Verb@sorted@ impose un ordre alphabétique stable
(\fichier{apiculture}, \fichier{app}, \fichier{camper}, \fichier{chasse}). C'est essentiel :
\fichier{vecteurs.npy} et \fichier{bm25.pkl} retrouvent les chunks \textbf{par position}, et
l'empreinte du cache d'embedding dépend de l'ordre.
\ligne{494--495} On ignore ce qui n'est pas un dossier (un fichier posé à la racine du corpus).
\ligne{496} Nom du dossier = domaine.
\ligne{497} Lecture du manifeste du domaine.
\ligne{498} \Verb@glob("*.md")@ : les fichiers Markdown du dossier, non récursif (les
sous-dossiers sont ignorés). \fichier{\_manifest.json} n'est pas concerné. Triés pour la même
raison qu'à la ligne 493.
\ligne{499} Ligne vide.
\ligne{500--501} Dossier sans \fichier{.md} : on passe au suivant.
\ligne{502--503} Affichage du nom du domaine entre crochets.
\end{explication}

\begin{shaded}
\begin{Verbatim}[numbers=left,firstnumber=504,numbersep=6pt,fontsize=\footnotesize,xleftmargin=1.2em]

        for f in fichiers:
            meta = {**config.META_DEFAUT, **manifeste.get(f.name, {})}
            # Si aucun manifeste, on derive un source_id du nom de fichier
            # pour que les identifiants de chunks restent uniques.
            if f.name not in manifeste:
                meta["source_id"] = slug(f.stem, 30)
                meta["source"] = f.stem.replace("_", " ")
            c = chunker_fichier(f, domaine, meta)
            tous.extend(c)
            if verbeux:
                print(f"    {f.name:38s} -> {len(c):3d} chunks")
\end{Verbatim}
\end{shaded}

\begin{explication}
\ligne{504} Ligne vide.
\ligne{505} Parcours des fichiers du domaine.
\ligne{506} \Verb@{**A, **B}@ \emph{fusionne} deux dictionnaires : on part de \Verb@META_DEFAUT@,
puis les clés du manifeste pour ce fichier \textbf{écrasent} les valeurs par défaut.
\Verb@manifeste.get(f.name, {})@ renvoie \Verb@{}@ si le fichier n'est pas décrit.
\ligne{507--508} Commentaires : sans manifeste, on dérive un \Verb@source_id@ du nom de fichier
pour que les identifiants restent uniques.
\ligne{509} Le fichier n'est pas décrit dans le manifeste…
\ligne{510} … \Verb@source_id@ = slug du nom sans extension (\Verb@f.stem@), limité à 30
caractères. Sans cela, tous ces fichiers auraient \Verb@source_id = "interne"@ et leurs
identifiants pourraient entrer en collision.
\ligne{511} … \Verb@source@ lisible : « aide\_application » $\rightarrow$ « aide application ».
\ligne{512} Découpage du fichier.
\ligne{513} \Verb@extend@ ajoute les \emph{éléments} de la liste \Verb@c@ à \Verb@tous@
(\Verb@append@ ajouterait la liste entière comme un seul élément).
\ligne{514--515} Affichage : \Verb@{f.name:38s}@ aligne le nom à gauche sur 38 caractères,
\Verb@{len(c):3d}@ écrit le nombre sur 3 caractères. On obtient une colonne bien alignée.
\end{explication}

\begin{shaded}
\begin{Verbatim}[numbers=left,firstnumber=516,numbersep=6pt,fontsize=\footnotesize,xleftmargin=1.2em]

    config.INDEX_DIR.mkdir(parents=True, exist_ok=True)
    config.CHUNKS_FILE.write_text(
        json.dumps(tous, ensure_ascii=False, indent=2), encoding="utf-8")

    if verbeux and tous:
        _rapport(tous)
    return tous
\end{Verbatim}
\end{shaded}

\begin{explication}
\ligne{516} Ligne vide.
\ligne{517} Création de \fichier{data/index/} : \Verb@parents=True@ crée aussi les dossiers
intermédiaires, \Verb@exist_ok=True@ évite une erreur si le dossier existe (équivalent de
\Verb@mkdir -p@).
\ligne{518--519} Écriture de \fichier{chunks.json}, qui écrase l'ancien fichier.
\Verb@ensure_ascii=False@ garde les accents lisibles (sinon « é » deviendrait
\Verb@é@) ; \Verb@indent=2@ rend le fichier lisible par un humain.
\ligne{520} Ligne vide.
\ligne{521--522} Rapport de contrôle seulement en mode verbeux \emph{et} s'il y a des chunks :
\Verb@_rapport@ appelle \Verb@min@ et divise par \Verb@len@, ce qui échouerait sur une liste vide.
\ligne{523} Renvoi de tous les chunks à l'appelant (\fichier{build\_embeddings.py}, qui les
passe ensuite à \Verb@embedder@).
\end{explication}

\subsection{\texorpdfstring{\texttt{\_rapport}}{\_rapport} et exécution directe (lignes 524 à 551)}
\begin{shaded}
\begin{Verbatim}[numbers=left,firstnumber=524,numbersep=6pt,fontsize=\footnotesize,xleftmargin=1.2em]


def _rapport(chunks: list[dict]) -> None:
    """Controle qualite affiche apres le chunking."""
    print(f"\n  TOTAL : {len(chunks)} chunks -> {config.CHUNKS_FILE}")

    for cle in ("domaine", "type"):
        compte = {}
        for c in chunks:
            v = c["metadata"][cle]
            compte[v] = compte.get(v, 0) + 1
        print(f"  Par {cle} : {compte}")
\end{Verbatim}
\end{shaded}

\begin{explication}
\ligne{524--525} Lignes vides.
\ligne{526} Définition : \Verb@-> None@, la fonction ne renvoie rien, elle affiche.
\ligne{527} Docstring : contrôle qualité affiché après le chunking.
\ligne{528} Nombre total de chunks et chemin du fichier.
\ligne{529} Ligne vide.
\ligne{530} Deux passages : un comptage par \Verb@domaine@, un par \Verb@type@.
\ligne{531} Dictionnaire de comptage \{valeur : nombre\}.
\ligne{532} Parcours de tous les chunks.
\ligne{533} Valeur de la métadonnée étudiée.
\ligne{534} Idiome de comptage : \Verb@compte.get(v, 0)@ renvoie 0 la première fois, puis on
ajoute 1. Équivalent de \Verb@collections.Counter@.
\ligne{535} Affichage, par exemple \newline{\footnotesize\Verb@Par domaine : {'apiculture': 61, 'app': 7, ...}@}\newline .
\end{explication}

\begin{shaded}
\begin{Verbatim}[numbers=left,firstnumber=536,numbersep=6pt,fontsize=\footnotesize,xleftmargin=1.2em]

    tailles = [c["metadata"]["nb_caracteres"] for c in chunks]
    print(f"  Taille : min={min(tailles)} / moy={sum(tailles)//len(tailles)} "
          f"/ max={max(tailles)}")

    courts = [c["id"] for c in chunks
              if c["metadata"]["nb_caracteres"] < config.TAILLE_MIN_ALERTE
              and c["metadata"]["type"] == "norme"]
    if courts:
        print(f"  Chunks courts (< {config.TAILLE_MIN_ALERTE} car.) : {len(courts)}"
              f" — sauves par le prefixe contextuel")
\end{Verbatim}
\end{shaded}

\begin{explication}
\ligne{536} Ligne vide.
\ligne{537} Liste des longueurs de tous les chunks.
\ligne{538--539} Minimum, moyenne et maximum. \Verb@//@ est la division \emph{entière}.
Deux f-strings consécutives sont concaténées automatiquement par Python.
\ligne{540} Ligne vide.
\ligne{541--543} Identifiants des chunks \Verb@"norme"@ de moins de 200 caractères.
\ligne{544} S'il y en a…
\ligne{545--546} … on affiche leur nombre, avec la mention « sauvés par le préfixe contextuel » :
un article très court reste compréhensible grâce à la source et au titre ajoutés.
\end{explication}

\begin{shaded}
\begin{Verbatim}[numbers=left,firstnumber=547,numbersep=6pt,fontsize=\footnotesize,xleftmargin=1.2em]


if __name__ == "__main__":
    print("\n=== ETAPE 1 : CHUNKING ===")
    chunker_corpus()
\end{Verbatim}
\end{shaded}

\begin{explication}
\ligne{547--548} Lignes vides.
\ligne{549} \Verb@__name__@ vaut \Verb@"__main__"@ seulement quand le fichier est
\emph{exécuté} (\Verb@python -m rag.chunker@), et non quand il est \emph{importé} par un autre
module. Le bloc suivant ne s'exécute donc pas lors d'un \Verb@import@.
\ligne{550} Affichage du titre de l'étape.
\ligne{551} Lancement du chunking complet, en mode verbeux par défaut. Dernière ligne du fichier.
\end{explication}


\newpage
% ============================================================================
\section{Remarques relevées pendant la lecture}
% ============================================================================

Rien de bloquant pour le fonctionnement actuel, sauf le point 1, qui a un effet réel.

\begin{enumerate}
  \item \textbf{Sections courtes perdues (lignes 385 et 390).} Une section qui n'est pas un
        article et fait 200 caractères ou moins est ignorée sans message. Dans
        \fichier{aide\_application.md}, trois réponses de la FAQ ne sont jamais indexées :
        « Puis-je créer une alerte sans réseau ? » (148 caractères), « Comment suivre mes
        signalements ? » (173) et « Les réponses du chatbot ont-elles une valeur légale ? »
        (160). Piste : un seuil plus bas pour les documents sans articles, ou aucun seuil
        quand le titre est une question.
  \item \textbf{\texttt{RE\_NOTE\_FIDELITE} n'est jamais utilisée (ligne 46).} Tout bloc
        \Verb@>@ qui n'est pas un avertissement est classé « note ».
  \item \textbf{Double point dans les chunks tableau (ligne 232)} : « (El khedra).. Date
        d'ouverture ».
  \item \textbf{Commentaire mal placé (ligne 53)} : il décrit la ligne 52.
  \item \textbf{Texte avant le premier titre perdu (ligne 129).}
  \item \textbf{Avertissements perdus hors article (lignes 376 et 390)} : aucun cas actuel.
  \item \textbf{« Article 1er » non reconnu (ligne 40)} à cause de \Verb@\b@ : le corpus
        actuel écrit « Article premier ».
  \item \textbf{Identifiants redondants (lignes 451--453)} pour les sections non-article, et
        collision possible si deux lignes de tableau ont les mêmes 30 premiers caractères.
  \item \textbf{\texttt{rang\_poids} n'est lu par aucun autre fichier (ligne 470).}
\end{enumerate}

% ============================================================================
\section{Résumé : le rôle de \texorpdfstring{\texttt{chunker.py}}{chunker.py} dans le RAG}
% ============================================================================

\fichier{chunker.py} est la \textbf{première étape} du RAG et elle en fixe la qualité : un
chunk mal découpé ne peut être rattrapé ni par le retriever, ni par le LLM. Le fichier
transforme les \fichier{.md} du corpus en environ 320 morceaux autonomes enregistrés dans
\fichier{data/index/chunks.json}.

\begin{itemize}
  \item \textbf{Découpage sémantique} : un article = un chunk, repéré par la regex
        « Article N », quel que soit le niveau de titre.
  \item \textbf{Contexte conservé} : une pile de parents garde la hiérarchie
        Titre $\rightarrow$ Chapitre $\rightarrow$ Article.
  \item \textbf{Cas particuliers} : tableaux découpés ligne par ligne, articles longs coupés
        par paragraphes (1500 caractères au maximum), article 12 découpé par gouvernorat,
        avertissements de concordance transformés en chunks à part.
  \item \textbf{Nettoyage} : notes éditoriales retirées ; coquilles officielles gardées, avec
        leur forme correcte en alias pour BM25.
  \item \textbf{Trois textes par chunk} : préfixe court + texte pour l'embedding, préfixe
        long + texte pour BM25, texte seul pour le LLM.
  \item \textbf{Métadonnées} (domaine, source, rang normatif, validité, article, type) : elles
        servent à \emph{filtrer} (spécialité de l'utilisateur), \emph{citer} (numéro d'article)
        et \emph{dater} (validité de la saison) dans les étapes suivantes.
\end{itemize}

Le fichier suivant dans le flux est \fichier{indexer.py} : il transforme ces chunks en
vecteurs (\fichier{vecteurs.npy}) et en index lexical (\fichier{bm25.pkl}).

\end{document}
